\documentclass[11pt]{article}

\usepackage[margin=1in]{geometry}
\usepackage{amsmath,amssymb,bm}
\usepackage{physics}
\usepackage{microtype}
\usepackage{hyperref}
\usepackage{enumitem}
\usepackage{booktabs}
\usepackage{mathtools}
\usepackage{array}
\usepackage{longtable}
\usepackage{amsthm}

\newtheorem{lemma}{Lemma}
\newtheorem{proposition}{Proposition}

\hypersetup{colorlinks=true,linkcolor=blue,citecolor=blue,urlcolor=blue}
\setlist{topsep=4pt,itemsep=2pt}

\title{\textbf{From Relational Progression to the Hypersurface-Deformation Algebra:\
A Conditional Route to Einstein Geometrodynamics and Universal Matter Cones}}
\author{Scott A. Jordan}
\date{August 28, 2026}

\begin{document}
\maketitle

\begin{abstract}
We use \emph{progression} as a term of art for locally realized physical evolution prior to assuming a spacetime time coordinate. A physical clock transition defines a pregeometric clock increment $d\sigma\equiv-(\hbar/\Delta E)d(\Delta\theta)$ and, relative to an arbitrary comparison parameter $\lambda$, a protocol-dependent representative $\Pi_{\mathcal C}=d\sigma/d\lambda$. Proper time is identified with $d\sigma$ only after geometric completion. We propose that many-fingered local progression defines an integrability problem on physical state space modulo spatial relabelings: the curvature of the local-update connection is required to be vertical, i.e. purely tangential to the relabeling orbit.

Assuming a smooth three-dimensional relational state manifold, a real linear extension of positive physical progression profiles to arbitrary canonical smearings, independence of normal and tangential generators, spatial covariance, and a leading bilinear single-gradient closure, the normal--normal commutator takes the form $[\delta_N,\delta_M]=\delta_{\beta^{ij}(N\partial_jM-M\partial_jN)}$. The Jacobi identities require $\beta^{ij}$ to be symmetric; positive nondegeneracy, pointwise locality, and a universal directional-calibration assumption then identify it with a spatial co-metric $h^{ij}$ and yield the undeformed hypersurface-deformation algebra (HDA).

Within the stated minimal local canonical class on $(h_{ij},\pi^{ij})$, closure fixes $b=-a/2$ and $ac=-1$ and excludes the tested $R^2$, $R_{ij}R^{ij}$, and $D^2R$ deformations. Scalar and Yang--Mills sectors share the same characteristic co-metric; the Dirac calculation establishes principal/characteristic matching only. The result is explicitly conditional: the proposed contribution is a possible progression-based physical origin of the HDA, not a new general proof that the HDA implies general relativity.\end{abstract}

\section{Introduction}

Quantum theory supplies a direct operational relation between energy and clock phase. In an already geometric relativistic description, a stationary energy component obeys
\begin{equation}
    d\theta=-\frac{E}{\hbar}\,d\tau,
    \label{eq:phase_proper}
\end{equation}
and for a physical clock built from two internal states separated by $\Delta E$, the observable relative phase satisfies
\begin{equation}
    d(\Delta\theta)=-\frac{\Delta E}{\hbar}\,d\tau.
    \label{eq:clock_phase}
\end{equation}
Matter-wave interferometry and quantum-clock interferometry exploit precisely this relation. Different spacetime histories generate different phases and, for genuine clock degrees of freedom, different accumulated proper times \cite{Stodolsky1979,Zych2011,Roura2020,Roura2021}. Relativistic atom interferometry provides a mature framework for connecting such phase evolution with gravitational motion and tests of universality \cite{Dimopoulos2008,Ufrecht2020}.

\paragraph{Terminology.}
Throughout this paper, \emph{progression} denotes locally realized physical evolution considered prior to committing to a spacetime time coordinate or external evolution parameter. The term does not denote a new substance, field, or second kind of time. To avoid importing proper time before geometry has been reconstructed, define the clock increment
\begin{equation}
    \boxed{d\sigma\equiv-\frac{\hbar}{\Delta E}\,d(\Delta\theta)}
    \label{eq:sigma_clock}
\end{equation}
for an ideal clock transition of fixed local splitting $\Delta E$. Before geometric completion, $\sigma$ is simply the operational reading accumulated along a relational clock history. Given an arbitrary comparison parameter $\lambda$, the quantity
\begin{equation}
    \Pi_{\mathcal C}\equiv\frac{d\sigma}{d\lambda}\bigg|_{\mathcal C}
\end{equation}
is a protocol- and parameter-dependent \emph{progression-rate representative}. After the relativistic geometry and common matter characteristic structure have been reconstructed, consistency with Eq.~\eqref{eq:clock_phase} identifies $d\sigma=d\tau$ for an ideal clock. Thus $d\sigma$ is the pregeometric operational increment, $\Pi_{\mathcal C}$ represents its rate relative to $\lambda$, $\lambda$ is only a comparison parameter, and the canonical lapse $N$ is a generator smearing. These objects may coincide numerically in an adapted representation but are not identified conceptually.

The present framework reverses the usual explanatory order. Rather than beginning with a spacetime metric and treating proper time as already given, it asks whether \emph{physical progression} can be treated as an operational primitive, revealed through local clock phase, while gravitational geometry emerges as the structure required for local relational evolution to be integrable and for universal conservation to remain self-consistent when progression varies from place to place. In this division of labor, quantum phase supplies an operational meaning for progression, integrability motivates the deformation structure, canonical closure fixes the minimal gravitational dynamics, and conservation supplies the covariant balance interpretation downstream.

Three objects must be distinguished from the outset. First is the \emph{progression principle}: local physical evolution is compared relationally rather than against a preferred external time. Second is an \emph{operational clock ratio}, defined only after specifying relational clock histories and a comparison protocol. Third is the canonical lapse $N$, an arbitrary smearing/refoliation variable in the Hamiltonian representation. The progression-rate representative $\Pi_{\mathcal C}$ may be assigned within a specified comparison protocol $\mathcal C$, but it is not itself an invariant local scalar independent of that protocol or of the chosen parameterization. The central hypothesis is instead that energetic inhomogeneity and local energy transfer are constrained by a universal progression-referenced conservation principle. In an adapted foliation the canonical lapse may represent the same local clock-rate assignment, but the paper does not identify the gauge variable $N$ ontologically with a material field $\Pi$.

Earlier formulations of this idea left two foundational arrows insufficiently justified. The first is more carefully stated here as
\begin{equation}
\begin{aligned}
    &dQ=\mathcal R_{\rm loc}\,d\sigma
      =\Pi_{\mathcal C}\mathcal R_{\rm loc}\,d\lambda
      \quad\text{(pregeometric temporal rate/measure conversion)}\\
    &\qquad\Longrightarrow\quad
      \nabla_\mu T^{\mu\nu}=0
      \ \text{with an emergent }g_{\mu\nu},
\end{aligned}
    \label{eq:first_arrow}
\end{equation}
and
\begin{equation}
\begin{aligned}
    &\text{energy of the progression/geometric response}\\
    &\qquad\Longrightarrow
    \text{ participation in the same total conservation balance}.
\end{aligned}
\label{eq:second_arrow}
\end{equation}
Neither follows from a single scalar current equation alone. The purpose of this paper is to replace those arrows with a sequence of explicit conditional steps and to identify exactly what remains to be established.

The main result is therefore best stated as a conditional reconstruction:
\begin{equation}
\boxed{
\begin{gathered}
\text{clock-defined relational progression on a smooth 3-manifold}\\
+\text{ local many-fingered evolution}\\
+\text{ vertical-curvature integrability modulo relabeling}\\
+\text{ real linear extension of positive-profile generators}\\
+\text{ independent normal/tangential generators}\\
+\text{ single-gradient locality}\\
+\text{ positive nondegenerate universal spatial calibration}\\
+\text{ the stated minimal canonical class}
\\[2pt]
\Longrightarrow\ \text{undeformed HDA}
\Longrightarrow\ \text{ADM/GR within that class}.
\end{gathered}}
\label{eq:conditional_reconstruction}
\end{equation}
The first implication is the progression-specific proposal; the second belongs to the established deformation-algebra/geometrodynamical tradition and is recalculated here in the restricted canonical class. The Tomonaga--Schwinger and Dirac--Schwinger traditions provide physical precedents for local-evolution integrability, but they presuppose causal/hypersurface structure and therefore do not derive the first implication. This distinction is the paper's intended novelty claim.

The resulting architecture has six layers. First, progression is defined relationally and operationally through the clock increment $d\sigma$, without using proper time as a primitive. Second, we introduce a local-evolution hypothesis: independently specifiable local progression acts as a many-fingered relational evolution. Motivated by Tomonaga--Schwinger integrability, different infinitesimal orderings that represent the same relational evolution are required to agree on physical observables, with spatial relabeling allowed as gauge. Under a minimal single-gradient closure this relational progression integrability yields a kinematic progression lemma whose structure tensor $\beta^{ij}$ is the candidate spatial co-metric; enforcing the Jacobi identities derives $\beta^{ij}=\beta^{ji}$ and constrains its intrinsic progression response to a symmetric tensor $K_{ij}$ under the minimal pointwise-local closure. Third, representing the resulting HDA canonically on $(h_{ij},\pi^{ij})$ selects the DeWitt kinetic combination and the spatial Ricci-scalar potential, and a broader four-spatial-derivative test eliminates the minimal $R^2$, $R_{ij}R^{ij}$, and $D^2R$ deformations. Fourth, parallel matter-closure calculations test a common characteristic geometry: scalar closure gives $\mathsf A^{CA}\mathsf C^{ij}_{AB}=\delta^C{}_B h^{ij}$, gauge closure gives $\alpha_E\alpha_B\chi^{ij}=h^{ij}$ modulo the Gauss constraint, and the Dirac principal/characteristic calculation gives $v^2\chi_D^{ij}=h^{ij}$, with the known standard tetrad--spinor branch supplying the lower-order Lorentz/gauge completion. Fifth, the $q=1$ statement is restricted pregeometrically to a one-factor temporal rate/measure conversion; after geometry emerges, an adapted lapse representative carries that factor inside $\sqrt{-g}=N\sqrt h$, while observer-current and Noether arguments give the covariant stress--energy interpretation. Sixth, the static two-function reduction connects the canonical completion to the previously derived weak-field quadratic progression stiffness, the Newtonian limit, exact spherical vacuum geometry, Maxwell propagation, and quantum phase-induced motion.

Throughout, we distinguish three kinds of statement: results derived within a stated mathematical completion; structural consequences of the progression premise; and hypotheses that remain to be derived. This distinction is essential because an Einstein--Hilbert completion reproduces general relativity in the sector where it is assumed. Any genuinely new phenomenology would have to arise from a controlled departure from that equilibrium completion rather than from reinterpretation alone.

\section{Operational progression from clock phase}

\subsection{Pregeometric clock increment and reparametrization invariance}

Let an ideal clock contain two internal states whose locally calibrated splitting is $\Delta E$. The directly registered quantity is their relative phase. We therefore take Eq.~\eqref{eq:sigma_clock}, rather than proper time, as the pregeometric operational definition of an infinitesimal clock increment. Introduce an arbitrary monotonic comparison parameter $\lambda$ along each relational clock history and define
\begin{equation}
    \boxed{d\sigma=\Pi_{\mathcal C}\,d\lambda.}
    \label{eq:progression_def}
\end{equation}
The parameter $\lambda$ is not assumed to be a universal physical time. The label $\mathcal C$ records the protocol by which separated clock histories are compared. Thus $\Pi_{\mathcal C}$ is a representative of a relational rate, not a pre-existing spacetime scalar.

For two specified clock histories $\Gamma$ and $\Gamma_0$ compared by the same protocol $\mathcal C$, define
\begin{equation}
    \boxed{
    \Pi_{\rm rel}[\Gamma,\Gamma_0;\mathcal C]
    =\frac{d\sigma_\Gamma/d\lambda}{d\sigma_{\Gamma_0}/d\lambda}.}
    \label{eq:pi_rel}
\end{equation}
Under $\lambda\mapsto f(\lambda)$ the common Jacobian cancels. Two abstract spatial labels alone therefore do not define an invariant clock ratio; a relational comparison protocol is part of the observable specification. After the spacetime geometry has been reconstructed, these histories become worldlines and the clock increment is identified with ordinary proper time for an ideal clock, $d\sigma=d\tau$. The standard reparametrization invariance of relativistic worldline actions then provides a downstream consistency check \cite{Roura2020}.

When the comparison protocol and an adapted representation are fixed, we suppress the label $\mathcal C$ and write $\Pi$. Nothing in Eq.~\eqref{eq:progression_def} identifies this representative with the ADM lapse as a physical field. In the canonical sections, $N$ is kept as an arbitrary external smearing variable; only after closure may an adapted foliation choose a particular $N$ that numerically represents the same clock-rate assignment. This separation is important because the off-shell hypersurface-deformation algebra is not simply an ontological statement about spacetime diffeomorphisms, and phase-space-dependent deformation parameters require additional care \cite{Thiemann2024,Bojowald2025}. It is then useful to introduce
\begin{equation}
    \boxed{A\equiv\ln\Pi.}
    \label{eq:A_def}
\end{equation}

\subsection{Clock phase and the later proper-time identification}

Equation~\eqref{eq:sigma_clock} immediately gives
\begin{equation}
    \Pi(x)=-\frac{\hbar}{\Delta E}\frac{\partial\Delta\theta}{\partial\lambda},
    \label{eq:pi_phase}
\end{equation}
and, for identical transitions,
\begin{equation}
    \boxed{
    \Pi_{\rm rel}[\Gamma,\Gamma_0;\mathcal C]
    =\frac{\partial_\lambda\Delta\theta_\Gamma}
           {\partial_\lambda\Delta\theta_{\Gamma_0}}.}
    \label{eq:pi_rel_phase}
\end{equation}
This definition uses an observable relative clock phase rather than a global energy-eigenstate phase \cite{Zych2011,Roura2020,Roura2021}. Once the emergent metric and common matter characteristic structure have been established, Eq.~\eqref{eq:clock_phase} gives $d\sigma=d\tau$. Only at that stage does the usual massive rest-energy phase take the geometric form
\begin{equation}
    d\theta_{\rm rest}=-\frac{mc^2}{\hbar}d\tau
    =-\frac{mc^2}{\hbar}d\sigma.
    \label{eq:rest_phase}
\end{equation}
Thus clock phase defines progression operationally upstream, while proper time and the rest-energy phase are recovered downstream as consistency relations of the completed geometry.

\section{From weighted progression to an invariant measure}

\subsection{Why one temporal rate/measure factor of \texorpdfstring{$\Pi$}{Pi} appears}

Consider a local transferred quantity $Q$ whose rate is measured per unit operational clock progression,
\begin{equation}
    \mathcal R_{\rm loc}=\frac{dQ}{d\sigma}.
\end{equation}
Using Eq.~\eqref{eq:progression_def},
\begin{equation}
    dQ
    =\mathcal R_{\rm loc}\,d\sigma
    =\Pi\,\mathcal R_{\rm loc}\,d\lambda,
    \qquad
    \frac{dQ}{d\lambda}=\Pi\mathcal R_{\rm loc}.
    \label{eq:rate_conversion}
\end{equation}
Thus one power of $\Pi$ is fixed when a local rate or temporal measure expressed per unit operational clock progression is rewritten per unit comparison parameter. We use $q=1$ as shorthand for this \emph{temporal rate/measure factor}.

No transformation law for an entire pregeometric four-current follows from Eq.~\eqref{eq:rate_conversion}. In particular, before a spatial flux structure and physical volume element have emerged, the statement
$J^\mu_{\lambda}=\Pi J^\mu_{\rm loc}$ is neither required nor assumed, and Eq.~\eqref{eq:rate_conversion} alone does not supply a four-dimensional divergence law. The complete local balance law is constructed only after spatial transport and the invariant measure are available. Conservation, equilibrium, and pointwise uniformity also remain distinct propositions; the one-factor rate conversion does not imply an algebraic relation such as $\rho\Pi=\text{constant}$ at each point.

\subsection{Including physical spatial volume}

A relativistic local conservation statement must refer not only to temporal rates but also to physical volume. Let a spatial slice carry metric $h_{ij}$ and volume element
\begin{equation}
    dV_{\rm phys}=\sqrt h\,d^3x.
\end{equation}
After a spatial geometry has emerged, the pregeometric clock increment is identified with the proper-time increment of an ideal clock, $d\sigma=d\tau$. A local matter action written in physical units then has the schematic form
\begin{equation}
    dS_m=\mathcal L_{\rm loc}\,d\sigma\,dV_{\rm phys}
        =\mathcal L_{\rm loc}\,d\tau\,dV_{\rm phys}.
\end{equation}
Using $d\sigma=\Pi d\lambda$ in an adapted representation,
\begin{equation}
    dS_m
    =\mathcal L_{\rm loc}\,\Pi\sqrt h\,d\lambda\,d^3x.
    \label{eq:measure_step}
\end{equation}
In an ADM decomposition, the invariant measure is
\begin{equation}
    \boxed{\sqrt{-g}=N\sqrt h.}
    \label{eq:measure_identity}
\end{equation}
For an adapted foliation and comparison protocol one may choose the coordinate lapse representative so that $N=\Pi_{\mathcal C}$ along the relevant congruence. In that representation, $\sqrt{-g}=\Pi_{\mathcal C}\sqrt h$. This is a gauge/representation statement, not an ontological identity between a canonical smearing variable and a physical scalar field. The same one-factor temporal conversion displayed in Eq.~\eqref{eq:rate_conversion} is then represented inside the invariant spacetime measure once spatial geometry is included.

This resolves an important bookkeeping issue. In a nongeometric description one may write an explicit factor $\Pi$. In a geometric description using $\sqrt{-g}$, that temporal measure factor is already present. Writing both without care would double-count the same conversion.

Once spatial transport has been defined and $J^\mu$ is a genuine spacetime current measured with respect to the emergent physical volume and proper progression, its densitized conservation law may be written in an adapted representation as
\begin{equation}
    \partial_\mu\!\left(\Pi\sqrt h\,J^\mu\right)=0,
\end{equation}
which, using Eq.~\eqref{eq:measure_identity}, is
\begin{equation}
    \boxed{
    \frac{1}{\sqrt{-g}}
    \partial_\mu\!\left(\sqrt{-g}\,J^\mu\right)
    =\nabla_\mu J^\mu=0.}
    \label{eq:cov_current}
\end{equation}
This is not yet a derivation of the full metric. It is a derivation of the form of the invariant measure \emph{conditional on} a physical spatial-volume structure and on choosing an adapted lapse representative for the operational progression assignment.

\subsection{What is and is not fixed by progression alone}

The progression factor fixes the temporal coframe in a chosen foliation,
\begin{equation}
    \vartheta^0=\Pi\,d\lambda.
\end{equation}
A complete local metric also requires spatial coframes $\vartheta^a$. The progression construction below does not by itself derive the relative Lorentzian sign. We select the Lorentzian/hyperbolic branch in accordance with the observed characteristic structure of relativistic matter and write
\begin{equation}
    g=-(\vartheta^0)^2+\delta_{ab}\vartheta^a\vartheta^b.
    \label{eq:coframe_metric}
\end{equation}
The signature choice is therefore an empirical/minimal-completion input at the present stage, not a consequence of progression integrability alone.
Thus $\Pi$ alone does not determine all of $g_{\mu\nu}$. In the static isotropic sector one scalar $B$ is sufficient to describe physical volume, but the fully dynamical theory also requires shift/vector and tensor degrees of freedom. The next section gives a kinematic route to this spatial structure: path independence of arbitrary local progression fixes the normal--normal update commutator up to a contravariant tensor $\beta^{ij}$; Jacobi closure derives its symmetry and constrains its intrinsic progression response, while directional flux calibration supplies the remaining positivity and nondegeneracy conditions needed to interpret $\beta^{ij}$ as $h^{ij}$.


\section{Relational progression integrability and the kinematic spatial tensor}
\label{sec:path_independence}

The progression premise can be strengthened in a way that does not yet assume a spacetime metric. Let $\Sigma$ be a smooth three-dimensional \emph{relational state manifold}: its points label simultaneously specified local physical data, but $\Sigma$ is not initially assumed to be embedded in a pre-existing spacetime. The existence, differentiable structure, and dimension of $\Sigma$ are assumptions of the present reconstruction; the framework derives a co-metric and deformation geometry on this manifold rather than deriving the manifold itself. Physical forward-progression assignments form a positive cone of scalar profiles,
\begin{equation}
    \Pi(x)>0,
\end{equation}
and a physical profile $\Pi$ generates an infinitesimal update $\delta_{\Pi}$.

The canonical commutator analysis requires a real linear space of test functions rather than only the positive cone. We therefore make an additional, explicit assumption: the physical update map is additive and positively homogeneous on the progression cone and admits a consistent local linear extension to arbitrary real smearings
\begin{equation}
    N,M\in C^\infty(\Sigma,\mathbb R).
    \label{eq:linear_extension_smearings}
\end{equation}
On any bounded local patch an arbitrary real smearing can be represented as a difference of positive profiles, and consistency of the extension makes the resulting generator independent of the chosen decomposition. We denote the extended generators by $\delta_N$ and $\delta_M$. Negative or sign-changing $N$ therefore need not be interpreted as negative physical progression; they are elements of the linearized generator space used to polarize and test the local algebra. The spatial generator $\delta_\xi$ denotes an infinitesimal relabeling generated by a vector field $\xi^i$ on $\Sigma$.

\subsection{State-space bundle formulation of ``the same relational evolution''}

To make the phrase ``the same relational evolution'' non-circular, let $\mathcal S$ denote the local state space over $\Sigma$ and let spatial relabelings act as a gauge group $\mathrm{Diff}(\Sigma)$. Physical instantaneous states are points of the quotient $\mathcal Q=\mathcal S/\mathrm{Diff}(\Sigma)$. Positive local progression profiles form a cone $\mathcal P_+$; the update map assigns to each profile a local vector field on $\mathcal S$. After the real linear extension described below, this assignment may be regarded as a connection-like horizontal lift from profile directions to state-space directions.

For two profiles $N$ and $M$, the commutator of their lifted updates is the infinitesimal curvature of this local-evolution connection. We define \emph{relational progression integrability modulo spatial gauge} to mean that this curvature is vertical with respect to the projection $\mathcal S\to\mathcal Q$: an infinitesimal update loop may move the representative state along a spatial-relabeling orbit but must induce no displacement in the physical quotient. Thus
\begin{equation}
    \boxed{\mathrm{pr}_{T\mathcal Q}[\,[\delta_N,\delta_M] \,]=0,}
    \label{eq:vertical_curvature}
\end{equation}
which is equivalently the statement that $[\delta_N,\delta_M]$ is a tangential/relabeling generator. This formulation supplies an independent mathematical meaning for ``the same relational evolution'': two paths are physically equivalent when they have the same endpoint in $\mathcal Q$, not because the desired deformation algebra has already been assumed. Global integrability may still fail through boundaries, topology, singularities, or genuine holonomy outside the local domain considered here.

\subsection{Physical basis: integrability of local relational evolution}

The additional physical hypothesis is that local progression can be specified independently in different regions, so that the state admits a many-fingered local evolution. This is not derived from clock phase alone. It is motivated by the Tomonaga--Schwinger formulation of relativistic quantum theory, in which a state functional on a spacelike hypersurface obeys a local evolution equation of the schematic form
\begin{equation}
    i\hbar\frac{\delta}{\delta\Sigma(x)}\ket{\Psi[\Sigma]}
    =\widehat{\mathcal H}(x)\ket{\Psi[\Sigma]},
    \label{eq:TS_equation}
\end{equation}
and consistency requires the result to be independent of the arbitrary ordering/foliation of compatible local deformations \cite{Tomonaga1946,Schwinger1948}. On a fixed spacetime this is closely tied to microcausality for spacelike-separated local Hamiltonian densities. The Dirac--Schwinger energy-density commutator and later integrability analyses show that, when local normal evolution is expressed in hypersurface language, its equal-surface commutator can contain a derivative-of-delta momentum-density term whose smeared form has the characteristic antisymmetric combination $N\partial_iM-M\partial_iN$ \cite{Schwinger1963,Fukuyama1973}. In canonical general relativity, Moncrief and Teitelboim showed conversely that the momentum constraints arise as integrability conditions of a Hamiltonian-constraint functional equation of Tomonaga--Schwinger type \cite{MoncriefTeitelboim1972}.

These results do not derive the present pregeometric construction: they already use spacelike hypersurfaces and causal or geometric structure. They instead supply a physical precedent for the principle we elevate one logical step earlier. We use the term \emph{relational progression integrability} for the vertical-curvature condition in Eq.~\eqref{eq:vertical_curvature}. Because points on the state surface have no physical identity beyond relational labels, a residual spatial relabeling is allowed. Thus integrability is required only on the quotient by spatial gauge, not as literal commutativity of the local normal generators.

There is also an operational clock reason to exclude a residual normal component. If an ordering loop generated an additional physical clock progression $\delta\sigma_{\rm hol}$, an ideal clock transition would acquire the measurable phase
\begin{equation}
    \delta(\Delta\theta)_{\rm hol}
    =-\frac{\Delta E}{\hbar}\,\delta\sigma_{\rm hol}.
    \label{eq:clock_holonomy}
\end{equation}
A residue that is only a spatial relabeling changes no relabeling-invariant clock observable, whereas a genuine residual normal progression would in principle leave an observable ordering-dependent clock phase. After geometric completion, ordinary proper-time differences between physically distinct worldlines remain allowed; what is excluded is phase/progression holonomy generated solely by an arbitrary decomposition or ordering of the same local relational evolution.

The requirement is first imposed on positive physical progression profiles and then extended bilinearly to the real test-smearing space of Eq.~\eqref{eq:linear_extension_smearings}. We therefore postulate
\begin{equation}
    \boxed{[\delta_N,\delta_M]=\delta_{\Xi(N,M)},}
    \label{eq:path_independence_postulate}
\end{equation}
where the right-hand side is purely tangential/spatial. We also retain the older descriptive phrase \emph{local progression-path consistency}; technically, Eq.~\eqref{eq:path_independence_postulate} states no residual normal holonomy after quotienting by spatial relabeling.

\subsection{Kinematic progression lemma}

\begin{lemma}[Kinematic progression lemma]
Let the local progression operators satisfy the following assumptions:
\begin{enumerate}[label=(\alph*)]
    \item \emph{Linear extension:} the physical positive-profile generators admit the real linear extension of Eq.~\eqref{eq:linear_extension_smearings}, so $\delta_{aN+bM}=a\delta_N+b\delta_M$ for the extended smearings.
    \item \emph{Relational progression integrability modulo labels:} Eq.~\eqref{eq:path_independence_postulate} holds, so different infinitesimal orderings of the same local relational evolution leave no additional physical progression after quotienting by spatial relabeling.
    \item \emph{Local single-gradient closure:} at a point $x$, $\Xi^i(N,M)$ is bilinear in $N,M$, depends only on their values and first derivatives at $x$, and contains at most one spatial derivative in total.
    \item \emph{Spatial covariance:} $N$ and $M$ are scalars under relabeling and $\Xi^i$ is a spatial vector.
    \item \emph{Generator independence:} normal-update and tangential-relabeling generators are independent directions in the local update algebra, so a vanishing sum of independent normal and tangential generators implies that their coefficients vanish separately.
\end{enumerate}
Then there exists a state-dependent contravariant tensor $\beta^{ij}(x)$, independent of the particular smearing functions $N$ and $M$, such that
\begin{equation}
    \boxed{
    [\delta_N,\delta_M]
    =\delta_{\xi},\qquad
    \xi^i=\beta^{ij}
    \left(N\partial_jM-M\partial_jN\right).}
    \label{eq:kinematic_lemma}
\end{equation}
\end{lemma}

\begin{proof}
Linearity of $\delta_N$ makes the commutator bilinear in $N$ and $M$, while antisymmetry of the commutator gives
\begin{equation}
    \Xi^i(N,M)=-\Xi^i(M,N).
\end{equation}
At a fixed point $p\in\Sigma$, locality and the single-gradient assumption allow the most general bilinear vector
\begin{equation}
    \Xi^i
    =A^i NM+B^{ij}N\partial_jM+C^{ij}M\partial_jN.
    \label{eq:general_Xi}
\end{equation}
Terms such as $D^{ijk}\partial_jN\partial_kM$ are excluded because they contain two spatial derivatives in total, while terms involving second derivatives are excluded by assumption (c). Antisymmetry under $N\leftrightarrow M$ first gives $A^i=0$ and then $C^{ij}=-B^{ij}$. Hence
\begin{equation}
    \Xi^i
    =B^{ij}\left(N\partial_jM-M\partial_jN\right).
\end{equation}
Defining $\beta^{ij}\equiv B^{ij}$ gives Eq.~\eqref{eq:kinematic_lemma}. Finally,
\begin{equation}
    \omega_j\equiv N\partial_jM-M\partial_jN
\end{equation}
is a spatial covector because $N$ and $M$ are scalars. Since $\Xi^i=\beta^{ij}\omega_j$ must transform as a vector for arbitrary $\omega_j$, $\beta^{ij}$ must transform as a rank-$(2,0)$ spatial tensor. The assumptions make it state dependent but independent of the particular pair $(N,M)$.
\end{proof}

The lemma is deliberately kinematic. By itself it proves neither that $\beta^{ij}$ is symmetric nor that it is the inverse of a positive spatial metric. The three-normal Jacobi identity will supply symmetry under the assumptions stated below; nondegeneracy, positivity, and the dynamical magnitude of its progression response require additional physical information.

\subsection{Local converse: the commutator implies integrability modulo gauge}

The physical comparison is between the two forward orderings
\begin{equation}
    e^{\epsilon\delta_N}e^{\eta\delta_M}
    \qquad\text{and}\qquad
    e^{\eta\delta_M}e^{\epsilon\delta_N}
\end{equation}
for positive physical profiles. Their leading mismatch is extracted algebraically by the closed product
\begin{equation}
    e^{\epsilon\delta_N}e^{\eta\delta_M}
    e^{-\epsilon\delta_N}e^{-\eta\delta_M}.
\end{equation}
The inverse exponentials here belong to the real linearized generator space used to diagnose the commutator; they do not represent a claim that physical clocks literally undergo backward progression. The Baker--Campbell--Hausdorff expansion gives
\begin{equation}
    1+\epsilon\eta[\delta_N,\delta_M]
    +O(\epsilon^2\eta,\epsilon\eta^2).
\end{equation}
If Eq.~\eqref{eq:kinematic_lemma} holds, the leading residue is only $\delta_\xi$. Therefore every observable invariant under spatial relabeling is unchanged by exchanging the local update ordering to this order. The result is a \emph{local} equivalence between the commutator relation and relational progression integrability modulo spatial gauge. Global integrability can still be obstructed by boundaries, topology, singularities, or nontrivial holonomy.

A useful check follows by taking $M=1$:
\begin{equation}
    [\delta_N,\delta_1]
    =-\delta_{\beta^{ij}\partial_jN}.
    \label{eq:uniform_check}
\end{equation}
Uniform progression commutes with another uniform progression, while a nonuniform progression profile differs in ordering from a uniform update by a purely spatial displacement proportional to its gradient. Likewise $M=cN$ gives a vanishing commutator, as expected for two proportional local progression profiles.

\subsection{Jacobi compatibility and the intrinsic progression response}
\label{sec:jacobi_compatibility}

The preceding commutator contains a state-dependent structure tensor. Its allowed evolution is therefore constrained by the Jacobi identities of the update operators. For this kinematic calculation write
\begin{equation}
    H[N]\equiv\delta_N,
    \qquad
    D[\xi]\equiv\delta_\xi,
\end{equation}
without yet assuming a Hamiltonian phase space. Define
\begin{equation}
    v^i(N,M)
    \equiv
    \beta^{ij}\omega_j(N,M),
    \qquad
    \omega_j(N,M)
    \equiv
    N\partial_jM-M\partial_jN,
    \label{eq:v_beta_omega}
\end{equation}
so that $[H[N],H[M]]=D[v(N,M)]$.

Because $v^i$ depends on the physical state through $\beta^{ij}$, the commutator with a state-dependent spatial generator contains an additional variation of its smearing vector:
\begin{equation}
    [H[N],D[\xi]]
    =D[\delta_N\xi]-H[\mathcal L_\xi N],
    \label{eq:state_dependent_HD}
\end{equation}
for the sign convention used in Eq.~\eqref{eq:tangent_normal} below. Applying Eq.~\eqref{eq:state_dependent_HD} to the three-normal Jacobi identity gives
\begin{equation}
\begin{split}
0={}&[H[N],[H[M],H[L]]]+\text{cyclic}\\
={}&D[\mathcal T]-H[\mathcal S],
\end{split}
\label{eq:three_normal_jacobi_split}
\end{equation}
where
\begin{align}
    \mathcal T^i
    &=(\delta_N\beta^{ij})\omega_j(M,L)
      +(\delta_M\beta^{ij})\omega_j(L,N)
      +(\delta_L\beta^{ij})\omega_j(N,M),
    \label{eq:jacobi_tangential}\\
    \mathcal S
    &=v^i(M,L)\partial_iN
      +v^i(L,N)\partial_iM
      +v^i(N,M)\partial_iL.
    \label{eq:jacobi_normal}
\end{align}
Assuming normal progression and spatial relabeling are independent directions in the update algebra, the two coefficients vanish separately.

\begin{proposition}[Jacobi symmetry condition]
The normal component $\mathcal S=0$ for arbitrary local smearings $N,M,L$ requires
\begin{equation}
    \boxed{\beta^{[ij]}=0.}
    \label{eq:beta_symmetric_jacobi}
\end{equation}
\end{proposition}

\begin{proof}
At a point $p$ choose $M=1$ and $N=L=0$ while keeping the gradients $\partial_iN=a_i$ and $\partial_iL=c_i$ arbitrary. Equation~\eqref{eq:jacobi_normal} becomes
\begin{equation}
    \mathcal S
    =\beta^{ij}(c_j a_i-a_j c_i)
    =2\beta^{[ij]}a_i c_j.
\end{equation}
Since $a_i$ and $c_j$ are arbitrary, $\mathcal S=0$ implies Eq.~\eqref{eq:beta_symmetric_jacobi}.
\end{proof}

Thus reciprocity of the spatial calibration is no longer required merely to obtain symmetry; Jacobi closure supplies the same condition kinematically. Reciprocity remains a useful physical interpretation of the resulting symmetric map.

The tangential component supplies the actual evolution compatibility condition:
\begin{equation}
\boxed{
(\delta_N\beta^{ij})\omega_j(M,L)
+(\delta_M\beta^{ij})\omega_j(L,N)
+(\delta_L\beta^{ij})\omega_j(N,M)=0.}
\label{eq:jacobi_beta_compatibility}
\end{equation}
Linearity of the progression generator implies that $\delta_N\beta^{ij}$ is linear in $N$. Under the same local first-gradient approximation used for the kinematic lemma, write
\begin{equation}
    \delta_N\beta^{ij}
    =NC^{ij}+C^{ijk}\partial_kN,
    \label{eq:beta_response_general}
\end{equation}
where the coefficients depend on the instantaneous state but not on the particular smearing function. Equation~\eqref{eq:beta_symmetric_jacobi} implies $C^{ij}=C^{ji}$ and $C^{ijk}=C^{jik}$.

The term proportional to $C^{ij}$ satisfies Eq.~\eqref{eq:jacobi_beta_compatibility} identically because
\begin{equation}
\begin{split}
&N\omega_j(M,L)+M\omega_j(L,N)+L\omega_j(N,M)=0.
\end{split}
\label{eq:cyclic_omega_identity}
\end{equation}
For the derivative term define
\begin{equation}
\begin{split}
Q_{jk}={}&(\partial_kN)\omega_j(M,L)
 +(\partial_kM)\omega_j(L,N)
 +(\partial_kL)\omega_j(N,M).
\end{split}
\label{eq:Q_jk}
\end{equation}
A direct interchange of $j$ and $k$ gives $Q_{jk}=-Q_{kj}$. Because arbitrary local smearings generate arbitrary antisymmetric two-covectors at a point, Jacobi requires
\begin{equation}
    C^{i[jk]}=0.
\end{equation}
Together with $C^{ijk}=C^{jik}$, this gives full symmetry of the rank-three coefficient:
\begin{equation}
    \boxed{
    \delta_N\beta^{ij}
    =NC^{(ij)}+C^{(ijk)}\partial_kN.}
    \label{eq:jacobi_allowed_beta_response}
\end{equation}

Equation~\eqref{eq:jacobi_allowed_beta_response} is the strongest consequence of the three-normal Jacobi identity at this derivative order. Jacobi does not force the intrinsic coefficient $C^{ij}$ to a particular value. A further progression-specific locality condition can simplify the result. If two profiles $N$ and $N'$ with the same value at $p$ produce the same intrinsic change of spatial calibration at that point,
\begin{equation}
    N(p)=N'(p)
    \quad\Longrightarrow\quad
    \delta_N\beta^{ij}(p)=\delta_{N'}\beta^{ij}(p),
    \label{eq:pointwise_progression_locality}
\end{equation}
then gradients of the smearing cannot enter the intrinsic response and
\begin{equation}
    \boxed{C^{(ijk)}=0,\qquad
    \delta_N\beta^{ij}=NC^{ij}.}
    \label{eq:minimal_beta_response}
\end{equation}
The same conclusion follows in a minimal tensor-content model in which there is no independent fully symmetric rank-three constitutive tensor available. Spatial gradients of $N$ still control the relative tilt/relabeling through Eq.~\eqref{eq:kinematic_lemma}; Eq.~\eqref{eq:minimal_beta_response} states only that the intrinsic pointwise calibration responds to the realized local progression itself.

If $\beta^{ij}$ is subsequently identified with a nondegenerate co-metric $h^{ij}$, differentiate $h_{ik}h^{kj}=\delta_i{}^j$ and define
\begin{equation}
    K_{ij}
    \equiv
    -\frac12 h_{ik}h_{j\ell}C^{k\ell}.
    \label{eq:K_from_C}
\end{equation}
Then the minimal Jacobi-compatible response becomes
\begin{equation}
    \boxed{
    \delta_N h_{ij}=2NK_{ij},
    \qquad
    \delta_N h^{ij}=-2NK^{ij}.}
    \label{eq:metric_progression_response}
\end{equation}
Including a spatial relabeling gives
\begin{equation}
    \delta_{(N,\xi)}h_{ij}
    =2NK_{ij}+\mathcal L_\xi h_{ij}.
    \label{eq:metric_progression_shift}
\end{equation}
The tensor $K_{ij}$ is an \emph{extrinsic-response tensor}: it measures how directional spatial calibration changes per unit realized local progression. Equation~\eqref{eq:metric_progression_response} has the familiar kinematic form of lapse-driven metric evolution, but here it is reached from the progression algebra rather than assumed as an embedding formula. The Jacobi identity does not determine $K_{ij}$ or its own evolution. That is the precise point at which kinematics ends and gravitational dynamics must begin.

\subsection{Mixed Jacobi identities}

The three-normal identity is the nontrivial condition that constrains the state-dependent structure tensor. The remaining local Jacobi identities provide consistency checks on the spatial-covariance assumptions. The $D,D,H$ identity follows from the representation property of Lie derivatives on the scalar smearing $N$. For the $D,H,H$ identity, use
\begin{equation}
    [D[\xi],H[N]]=H[\mathcal L_\xi N],
    \qquad
    [H[N],H[M]]=D[v(N,M)],
\end{equation}
with $v^i=\beta^{ij}(N\partial_jM-M\partial_jN)$. If $\beta^{ij}$ transforms as a spatial $(2,0)$ tensor, direct substitution gives
\begin{equation}
    \mathcal L_\xi v(N,M)
    =v(\mathcal L_\xi N,M)+v(N,\mathcal L_\xi M)
      +(\mathcal L_\xi\beta)^{ij}\omega_j(N,M),
\end{equation}
and the last term is precisely the variation of the state-dependent structure tensor under the relabeling generator. Hence the mixed Jacobi identity vanishes. The complete kinematic algebra is therefore Jacobi compatible at this derivative order provided the spatial generator acts tensorially on $\beta^{ij}$ and the normal/tangential independence assumption used above holds.

\subsection{From \texorpdfstring{$\beta^{ij}$}{beta} to a spatial co-metric}

The tensor $\beta^{ij}$ can acquire an operational calibration from directional momentum transport. Let $s_i$ be the normal covector to a small oriented surface element. Local momentum-flux measurements require some contravariant map from surface orientation to transport direction. In the minimal universal completion we make the additional identification that this transport map is the same state-space tensor already selected by the progression commutator,
\begin{equation}
    s_i\longmapsto s^i=\beta^{ij}s_j.
    \label{eq:direction_map}
\end{equation}
This identification is a universality/calibration assumption, not a consequence of the Jacobi identity itself.
The Jacobi analysis in Sec.~\ref{sec:jacobi_compatibility} already enforces $\beta^{ij}=\beta^{ji}$ under the stated independence assumptions. Reciprocity of directional transport supplies an independent physical interpretation of that symmetry. If the symmetric map is also nondegenerate and positive on spatial directions,
\begin{equation}
    \beta^{ij}s_i s_j>0
    \qquad(s_i\neq0),
\end{equation}
then $\beta^{ij}$ possesses a positive-definite inverse. Define
\begin{equation}
    \boxed{h_{ij}\equiv(\beta^{-1})_{ij},\qquad \beta^{ij}=h^{ij}.}
    \label{eq:beta_metric}
\end{equation}
Thus the first genuinely spatial tensor can, under this universal calibration assumption, be interpreted as the reciprocal directional map used by local momentum transport. Symmetry follows from Jacobi compatibility once the progression and relabeling generators are treated as independent; nondegeneracy and positivity remain additional physical requirements that must be justified by the microscopic progression theory.

Once $h_{ij}$ is available, compatibility of the finite-region conservation measure with the same spatial structure gives the minimal choice
\begin{equation}
    dV_{\rm phys}=\sqrt h\,d^3x.
\end{equation}
An additional independent factor $f(x)\sqrt h$ would introduce a second scalar spatial density not fixed by the progression construction. Excluding such an extra degree of freedom is therefore part of the minimal completion, not a theorem of path independence alone.

\subsection{Completion to the hypersurface-deformation algebra}

Spatial relabelings close under the ordinary Lie bracket,
\begin{equation}
    [\delta_\xi,\delta_\eta]
    =\delta_{[\xi,\eta]},
    \qquad
    [\xi,\eta]^i
    =\xi^j\partial_j\eta^i-\eta^j\partial_j\xi^i.
    \label{eq:tangent_tangent}
\end{equation}
Because $N$ is a spatial scalar,
\begin{equation}
    [\delta_\xi,\delta_N]
    =\delta_{\mathcal L_\xi N},
    \label{eq:tangent_normal}
\end{equation}
up to the conventional overall sign assigned to the generators. Combining Eqs.~\eqref{eq:kinematic_lemma}, \eqref{eq:beta_metric}, \eqref{eq:tangent_tangent}, and \eqref{eq:tangent_normal} gives
\begin{equation}
\boxed{
\begin{aligned}
[\delta_\xi,\delta_\eta]
&=\delta_{[\xi,\eta]},\\
[\delta_\xi,\delta_N]
&=\delta_{\mathcal L_\xi N},\\
[\delta_N,\delta_M]
&=\delta_{h^{ij}(N\partial_jM-M\partial_jN)}.
\end{aligned}}
\label{eq:HDA_kinematic}
\end{equation}
This is the kinematic hypersurface-deformation structure. Because $h^{ij}$ depends on the state, Eq.~\eqref{eq:HDA_kinematic} is more precisely a Lie-algebroid-type relation than an ordinary Lie algebra with constant structure constants.

Teitelboim derived the corresponding constraint commutators from path independence of evolution between hypersurfaces, and Hojman, Kucha\v{r}, and Teitelboim subsequently asked what canonical geometrodynamics realizes that deformation structure \cite{Teitelboim1973,HKT1976}. The progression-specific proposal is logically earlier: derive the normal-deformation structure itself from arbitrary local progression and conservation, rather than beginning by assuming a spacetime embedding. This distinction is also the principal point at which the proposed derivation remains incomplete.

\subsection{Canonical representation}

If the physical state space admits a Hamiltonian description, let $H[N]$ generate local progression and $D[\xi]$ generate spatial relabeling. For a phase-space functional $F$,
\begin{equation}
    \delta_{N,\xi}F
    =\left\{F,H[N]+D[\xi]\right\}.
\end{equation}
Hamiltonian vector fields satisfy $[X_A,X_B]=X_{\{A,B\}}$, so Eq.~\eqref{eq:HDA_kinematic} is represented by
\begin{equation}
\boxed{
\begin{aligned}
\{D[\xi],D[\eta]\}
&=D[[\xi,\eta]],\\
\{D[\xi],H[N]\}
&=H[\mathcal L_\xi N],\\
\{H[N],H[M]\}
&=D[h^{ij}(N\partial_jM-M\partial_jN)].
\end{aligned}}
\label{eq:HDA_PB}
\end{equation}
Hojman, Kucha\v{r}, and Teitelboim showed that representing this deformation structure under suitable canonical assumptions strongly restricts geometrodynamics \cite{HKT1976}. Here we perform the minimal bracket calculation explicitly rather than leaving canonical closure as a future invocation. The result is developed in the next section.

A technical separation will be maintained throughout. The functions $N$ and $N^i$ that smear the canonical constraints are treated as arbitrary external deformation parameters when the Poisson algebra is calculated. Only after closure has been established may one select an adapted foliation in which a particular lapse profile represents an operational clock/progression assignment. This avoids extra bracket terms that can arise if lapse or shift are made phase-space dependent during the algebraic derivation and keeps the off-shell HDA conceptually distinct from an ontological scalar progression field \cite{Thiemann2024,Bojowald2025}.

The kinematic chain established before canonical closure is
\begin{equation}
\boxed{
\begin{aligned}
&\text{arbitrary local progression}
+\text{local-evolution hypothesis}\\
&\quad+\text{relational progression integrability modulo labels}\\
&\quad+\text{local single-gradient closure}
+\text{spatial covariance}
\Longrightarrow \beta^{ij},\\
&\text{three-normal Jacobi closure}
\Longrightarrow \beta^{ij}=\beta^{ji}
\text{ and Eq.~\eqref{eq:jacobi_allowed_beta_response}},\\
&\text{nondegenerate positive calibration}
\Longrightarrow \beta^{ij}=h^{ij},\\
&\text{pointwise progression locality}
\Longrightarrow \delta_Nh_{ij}=2NK_{ij},\\
&\text{spatial relabeling covariance}
\Longrightarrow \text{kinematic HDA}.
\end{aligned}}
\label{eq:path_to_HDA}
\end{equation}
The remaining question at this stage is whether a minimally assumed canonical generator is fixed by Eq.~\eqref{eq:HDA_PB} strongly enough to determine $K_{ij}$ and the spatial potential.

\subsection{Scope and failure modes of the lemma}

Four assumptions now carry most of the foundational burden. First, relational progression integrability remains a physical hypothesis rather than a consequence of the arbitrariness of $\lambda$. Tomonaga--Schwinger theory motivates integrability of local relativistic evolution, and Eq.~\eqref{eq:clock_holonomy} gives an operational reason to exclude ordering-dependent normal progression, but neither result proves that a pregeometric progression theory admits the required many-fingered local evolution or only spatial-gauge holonomy. If equivalent local update orderings leave a non-gauge residual normal update, Eq.~\eqref{eq:path_independence_postulate} is too strong. This condition does not forbid ordinary proper-time differences between physically distinct worldlines. Second, the single-gradient assumption excludes otherwise admissible antisymmetric terms such as $D^{ijk}\partial_jN\partial_kM$ and all higher-derivative structures. If such terms are physically required, the progression kinematics need not reduce to the standard hypersurface-deformation algebra. Third, although Jacobi closure derives symmetry of $\beta^{ij}$, identifying it with an ordinary spatial co-metric still requires nondegeneracy and positivity. Fourth, the minimal intrinsic law $\delta_N\beta^{ij}=NC^{ij}$ requires the additional pointwise-progression-locality condition in Eq.~\eqref{eq:pointwise_progression_locality}, or an equivalent restriction on tensor content. If the permitted $C^{(ijk)}\partial_kN$ term is physical, the geometric interpretation must be enlarged rather than silently discarded.


\section{Canonical progression closure and emergent gravitational dynamics}
\label{sec:canonical_closure}

The kinematic analysis determines the deformation algebra but, by itself, does not determine the symmetric response tensor $K_{ij}$. We now ask whether the simplest local canonical representation of the progression HDA fixes that freedom. The result is conditional on a minimal phase space and locality assumptions, but it can be calculated directly without starting from the Einstein--Hilbert action.

Throughout this section, $R\equiv{}^{(3)}R$ denotes the scalar curvature of $h_{ij}$.

\subsection{General quadratic local gravitational constraint within the stated class}

Take the canonical pair
\begin{equation}
    \{h_{ij}(x),\pi^{kl}(y)\}
    =\delta_{(i}^{k}\delta_{j)}^{l}\,\delta^3(x-y),
    \label{eq:canonical_pair}
\end{equation}
and the spatial relabeling generator
\begin{equation}
    D[\xi]
    =\int d^3x\,\pi^{ij}\mathcal L_\xi h_{ij}
    =-2\int d^3x\,\xi_jD_i\pi^{ij},
    \label{eq:D_generator}
\end{equation}
where boundary terms are suppressed. Consider the most general ultralocal quadratic momentum term and the lowest-derivative scalar potential,
\begin{equation}
\boxed{
\mathcal H_\perp
=\frac{1}{\sqrt h}
\left(a\pi^{ij}\pi_{ij}+b\pi^2\right)
+\sqrt h\left(cR+d\right),}
\label{eq:abcd_constraint}
\end{equation}
with
\begin{equation}
    H[N]=\int d^3x\,N\mathcal H_\perp,
    \qquad
    \pi=h_{ij}\pi^{ij}.
\end{equation}
At this stage $a,b,c,d$ are independent constants. Define
\begin{equation}
    w_i\equiv ND_iM-MD_iN.
    \label{eq:w_def}
\end{equation}
The required progression HDA is
\begin{equation}
    \{H[N],H[M]\}
    =D[h^{ij}w_j].
    \label{eq:required_HDA}
\end{equation}

Because the kinetic term is ultralocal and the potential contains no momenta, the antisymmetrized kinetic--kinetic and potential--potential brackets vanish. The only nontrivial contribution is the cross bracket. The momentum derivative of the kinetic term is
\begin{equation}
    \frac{\delta T[N]}{\delta\pi^{ij}}
    =\frac{2N}{\sqrt h}
    \left(a\pi_{ij}+b\pi h_{ij}\right).
    \label{eq:dT_dpi}
\end{equation}
The derivative-of-smearing part of the curvature variation is
\begin{equation}
    \left.
    \frac{\delta V[M]}{\delta h_{ij}}
    \right|_{\partial M}
    =c\sqrt h\left(D^iD^jM-h^{ij}D^2M\right).
    \label{eq:dV_dh}
\end{equation}
Terms proportional to $M$ without derivatives cancel under $N\leftrightarrow M$, including the constant $d$ term. A direct contraction and one integration by parts give
\begin{equation}
\boxed{
\begin{aligned}
\{H[N],H[M]\}
={}&-ac\,D[h^{ij}w_j]\\
&-2c(a+2b)\int d^3x\,w^iD_i\pi .
\end{aligned}}
\label{eq:abcd_bracket}
\end{equation}
Exact progression-HDA closure on unrestricted phase space therefore requires
\begin{equation}
    \boxed{b=-\frac a2,\qquad ac=-1.}
    \label{eq:abcd_conditions}
\end{equation}
The constant $d$ is not fixed by the bracket. Writing
\begin{equation}
    d=\frac{2\Lambda}{a},
\end{equation}
the constraint becomes
\begin{equation}
\boxed{
\mathcal H_\perp
=\frac{a}{\sqrt h}
\left(\pi^{ij}\pi_{ij}-\frac12\pi^2\right)
-\frac{\sqrt h}{a}
\left(R-2\Lambda\right).}
\label{eq:ADM_selected_constraint}
\end{equation}
Thus the HDA selects the DeWitt trace coefficient and the relative kinetic/curvature normalization within the ansatz. Equivalently, if the trace term is parameterized as $-\lambda_{\rm tr}\pi^2$, unrestricted closure selects $\lambda_{\rm tr}=1/2$. Values $\lambda_{\rm tr}\neq1/2$ require an additional condition such as $D_i\pi=0$ or a modified constraint structure and therefore do not represent the undeformed progression HDA on the full phase space.

\subsection{Canonical determination of the extrinsic response}

Hamilton's equation generated by Eq.~\eqref{eq:ADM_selected_constraint} is
\begin{equation}
\boxed{
\delta_Nh_{ij}
=\frac{2aN}{\sqrt h}
\left(\pi_{ij}-\frac12\pi h_{ij}\right).}
\label{eq:h_evolution_canonical}
\end{equation}
Comparing with the Jacobi-compatible kinematic form $\delta_Nh_{ij}=2NK_{ij}$ gives
\begin{equation}
\boxed{
K_{ij}
=\frac{a}{\sqrt h}
\left(\pi_{ij}-\frac12\pi h_{ij}\right),}
\label{eq:K_pi_relation}
\end{equation}
and inversion yields
\begin{equation}
\boxed{
\pi^{ij}
=\frac{\sqrt h}{a}
\left(K^{ij}-Kh^{ij}\right).}
\label{eq:pi_K_relation}
\end{equation}
The symmetric tensor left free by the Jacobi identity is therefore determined canonically once the minimal HDA representation is required.

Including the shift, the canonical action is
\begin{equation}
    S=\int d\lambda\,d^3x
    \left(\pi^{ij}\dot h_{ij}
    -N\mathcal H_\perp-N^i\mathcal H_i\right).
\end{equation}
Eliminating $\pi^{ij}$ with Eq.~\eqref{eq:pi_K_relation} gives, up to the standard boundary term,
\begin{equation}
\boxed{
S_g
=\frac{1}{a}\int d\lambda\,d^3x\,
N\sqrt h\left(
K_{ij}K^{ij}-K^2+R-2\Lambda
\right).}
\label{eq:ADM_reconstructed}
\end{equation}
The overall coefficient $1/a$ is not fixed by the algebra; it is the gravitational coupling normalization. With conventional units it is fixed by the measured Newtonian coupling and reproduces the usual ADM/Einstein--Hilbert normalization.

\subsection{Relaxing the spatial potential through four derivatives}

Equation~\eqref{eq:abcd_constraint} still began with an Einstein-shaped two-derivative potential. To test whether that shape was doing hidden work, retain the already selected DeWitt kinetic term but enlarge the parity-even local spatial potential through four derivatives. In three dimensions a convenient independent basis is
\begin{equation}
\boxed{
U[h]
=d+cR+\alpha R^2+\zeta R_{ij}R^{ij}+\gamma D^2R,}
\label{eq:four_derivative_potential}
\end{equation}
where an $R_{ijkl}R^{ijkl}$ term is redundant because
\begin{equation}
    R_{ijkl}R^{ijkl}=4R_{ij}R^{ij}-R^2
\end{equation}
in three dimensions. Parity-odd densities and terms with six or more spatial derivatives are not included in this test.

It is useful first to treat the entire $f(R)$ class. Let
\begin{equation}
    V_f[N]=\int d^3x\,N\sqrt h\,f(R),
    \qquad F(R)=f'(R).
\end{equation}
With
\begin{equation}
    P_{ij}\equiv\pi_{ij}-\frac12\pi h_{ij},
\end{equation}
the derivative-of-smearing variation is proportional to
\begin{equation}
    D^iD^j(MF)-h^{ij}D^2(MF),
\end{equation}
and the DeWitt contraction removes the trace piece. The exact antisymmetrized bracket contribution is
\begin{equation}
\boxed{
\{H[N],H[M]\}_{f}
=2a\int d^3x\,
F\,w_jD_i\pi^{ij}
-2a\int d^3x\,
w_j\pi^{ij}D_iF.}
\label{eq:fR_bracket}
\end{equation}
The second term has no counterpart in the progression HDA. For arbitrary canonical data it requires $D_iF=0$, and therefore
\begin{equation}
    \boxed{f''(R)=0,\qquad f(R)=cR+d.}
    \label{eq:fR_linear}
\end{equation}
In particular, an $R^2$ correction gives the anomaly
\begin{equation}
\boxed{
\Delta_{R^2}
=4a\alpha\int d^3x\,w_j
\left(RD_i\pi^{ij}-\pi^{ij}D_iR\right),}
\label{eq:R2_anomaly}
\end{equation}
so unrestricted closure requires $\alpha=0$.

For the Ricci-squared term define $Q=R_{ij}R^{ij}$. The derivative-dependent part of the metric variation may be written
\begin{equation}
\begin{aligned}
\mathcal Q^{ij}[M]
={}&D_kD^i(MR^{kj})+D_kD^j(MR^{ki})\\
&-D^2(MR^{ij})-h^{ij}D_kD_l(MR^{kl}).
\end{aligned}
\label{eq:Q_variation}
\end{equation}
Its principal symbol is
\begin{equation}
\boxed{
A^{ij\,kl}
=h^{i(k}R^{l)j}+h^{j(k}R^{l)i}
-R^{ij}h^{kl}-h^{ij}R^{kl}.}
\label{eq:Ricci_principal_symbol}
\end{equation}
At a point where derivatives of the curvature are set to zero for the principal-symbol test, the bracket contains
\begin{equation}
    \Delta_{R_{ij}^2}^{\rm pr}
    =2a\zeta\int d^3x\,
    w_l A^{ij\,kl}D_kP_{ij}+\cdots .
    \label{eq:Ricci2_anomaly}
\end{equation}
The obstruction can be exhibited explicitly. Choose Riemann-normal spatial coordinates at the test point so that $h_{ij}=\delta_{ij}$ and take
\begin{equation}
    R_{ij}=\operatorname{diag}(r_1,r_2,r_3),
    \qquad r_2\neq r_3.
\end{equation}
Let the only nonzero first derivatives of $P_{ij}$ at that point be
\begin{equation}
    D_1P_{22}=X,
    \qquad
    D_1P_{33}=-X.
\end{equation}
Then $D_kP=0$ and $D_iP^{ij}=0$, so the momentum-constraint structure that could appear on the right-hand side of the undeformed HDA vanishes. By contrast, Eq.~\eqref{eq:Ricci_principal_symbol} gives
\begin{equation}
    A^{22\,11}=-(r_1+r_2),
    \qquad
    A^{33\,11}=-(r_1+r_3),
\end{equation}
and therefore
\begin{equation}
    A^{ij\,k1}D_kP_{ij}=X(r_3-r_2)\neq0.
\end{equation}
For a lapse pair with $w_1\neq0$, the Ricci-squared principal contribution is consequently nonzero even though the required momentum-constraint term is zero. No nonzero constant $\zeta$ can reproduce the $h^{ij}$ structure for arbitrary data. Hence
\begin{equation}
    \boxed{\zeta=0.}
    \label{eq:zeta_zero}
\end{equation}
Finally, the $D^2R$ term can be integrated by parts as $\int\sqrt h\,R D^2N$ up to a boundary term. Its curvature variation therefore produces a highest-derivative lapse contribution
\begin{equation}
\boxed{
\Delta_{D^2R}^{\rm high}
=-2a\gamma\int d^3x\,\pi^{ij}
\left[ND_iD_jD^2M-MD_iD_jD^2N\right]+\cdots .}
\label{eq:D2R_anomaly}
\end{equation}
The undeformed progression HDA is first order in the lapse profiles and contains no structure able to cancel these fourth derivatives. Thus
\begin{equation}
    \boxed{\gamma=0.}
\end{equation}
Within the parity-even four-spatial-derivative basis, exact progression-HDA closure therefore selects
\begin{equation}
\boxed{
U[h]= -\frac{1}{a}\left(R-2\Lambda\right),}
\label{eq:potential_uniqueness}
\end{equation}
with the overall gravitational coupling and $\Lambda$ left free. This is a conditional uniqueness statement: higher-curvature generally covariant theories may require enlarged phase spaces, while preferred-foliation theories may consistently realize a different constraint algebra. The result applies to the minimal $(h_{ij},\pi^{ij})$ phase space with the undeformed progression HDA.

\subsection{Scalar matter HDA closure}
\label{subsec:matter_closure}

The same closure condition can be imposed on matter rather than appending universal metric coupling afterward. Consider scalar fields $\phi^A$ with conjugate momenta $p_A$,
\begin{equation}
    \{\phi^A(x),p_B(y)\}=\delta^A_B\delta^3(x-y),
\end{equation}
and the broad minimally separable local Hamiltonian
\begin{equation}
\boxed{
\begin{aligned}
\mathcal H_m={}&
\frac{1}{2\sqrt h}
 p_A\mathsf A^{AB}(\phi)p_B\\
&+\frac{\sqrt h}{2}
\mathsf C^{ij}_{AB}(\phi,h)
D_i\phi^A D_j\phi^B
+\sqrt h\,U(\phi).
\end{aligned}}
\label{eq:general_matter_H}
\end{equation}
The matter momentum generator required by spatial relabeling covariance is
\begin{equation}
    \boxed{D_m[\xi]=\int d^3x\,\xi^i p_A D_i\phi^A.}
    \label{eq:D_matter}
\end{equation}
The momentum derivative is
\begin{equation}
    \frac{\delta H_m[N]}{\delta p_A}
    =\frac{N}{\sqrt h}\mathsf A^{AB}p_B,
\end{equation}
and the only part of the field derivative relevant after antisymmetrizing the lapse profiles is
\begin{equation}
    \left.
    \frac{\delta H_m[N]}{\delta\phi^A}
    \right|_{\partial N}
    =-D_i\left(N\sqrt h\,\mathsf C^{ij}_{AB}D_j\phi^B\right).
\end{equation}
Therefore
\begin{equation}
\boxed{
\{H_m[N],H_m[M]\}
=\int d^3x\,w_i\,
 p_C\mathsf A^{CA}\mathsf C^{ij}_{AB}D_j\phi^B.}
\label{eq:matter_bracket}
\end{equation}
Exact closure with the same progression structure function requires
\begin{equation}
    \{H_m[N],H_m[M]\}=D_m[h^{ij}w_j],
\end{equation}
and hence, for arbitrary matter data,
\begin{equation}
\boxed{
\mathsf A^{CA}\mathsf C^{ij}_{AB}
=\delta^C{}_B h^{ij}.}
\label{eq:matter_closure_condition}
\end{equation}
If $\mathsf A$ is invertible,
\begin{equation}
\boxed{
\mathsf C^{ij}_{AB}
=(\mathsf A^{-1})_{AB}h^{ij}.}
\label{eq:C_metric_factorization}
\end{equation}
Thus the same co-metric selected by progression kinematics governs the principal spatial propagation of all fields in this minimal scalar class. Field-space normalizations may differ, and $U(\phi)$ is unconstrained by the HDA because it produces no derivative-of-lapse terms, but there is no independent species-dependent spatial propagation metric.

For the minimal coupling in Eq.~\eqref{eq:general_matter_H}, $H_m$ depends algebraically on $h_{ij}$ and not on $\pi^{ij}$ or derivatives of $h_{ij}$. The gravity--matter cross brackets are therefore proportional to $NM$ and cancel under antisymmetrization. Consequently
\begin{equation}
\boxed{
\{H_{\rm tot}[N],H_{\rm tot}[M]\}
=D_{\rm tot}[h^{ij}w_j].}
\label{eq:total_HDA_matter}
\end{equation}
The geometric and matter sectors form one first-class deformation system.

The scalar calculation already fixes the physical co-metric for a broad second-order matter prototype. Gauge fields and spinors have additional internal constraints, so their closure must be checked separately rather than inferred from the scalar result.

\subsection{Maxwell and Yang--Mills closure modulo the Gauss constraint}
\label{subsec:YM_closure}

Let $A_i^a$ be a spatial gauge connection and $E^{ia}$ its conjugate electric momentum,
\begin{equation}
    \{A_i^a(x),E^{jb}(y)\}
    =\delta_i^j\delta^{ab}\delta^3(x-y),
\end{equation}
with curvature
\begin{equation}
    F_{ij}^a=2\partial_{[i}A_{j]}^a+f^a{}_{bc}A_i^bA_j^c.
\end{equation}
The Maxwell case is obtained by suppressing the internal index and setting $f^a{}_{bc}=0$. Let $\kappa_{ab}$ be an invariant positive inner product on the gauge algebra. To test whether the gauge sector could carry an independent propagation geometry, introduce a candidate co-metric $\chi^{ij}$ with inverse $\chi_{ij}$ and write the local trial Hamiltonian
\begin{equation}
\boxed{
H_{\rm YM}[N]
=\int d^3x\,N\left[
\frac{\alpha_E}{2\sqrt h}\chi_{ij}\kappa_{ab}E^{ia}E^{jb}
+\frac{\alpha_B\sqrt h}{4}\chi^{ik}\chi^{jl}
\kappa_{ab}F^a_{ij}F^b_{kl}
\right].}
\label{eq:YM_trial_H}
\end{equation}
The canonical derivatives relevant to the antisymmetrized lapse bracket are
\begin{equation}
    \frac{\delta H_{\rm YM}[N]}{\delta E^{ia}}
    =\frac{\alpha_E N}{\sqrt h}\chi_{ij}\kappa_{ab}E^{jb},
\end{equation}
\begin{equation}
    \frac{\delta H_{\rm YM}[N]}{\delta A_i^a}
    =-\alpha_B\mathcal D_j\!\left(
    N\sqrt h\,\chi^{jk}\chi^{il}\kappa_{ab}F^b_{kl}
    \right),
\end{equation}
where $\mathcal D_i$ is the gauge-covariant derivative. Terms with no derivative of $N$ or $M$ cancel after antisymmetrization. Direct substitution therefore gives
\begin{equation}
\boxed{
\{H_{\rm YM}[N],H_{\rm YM}[M]\}
=V_{\rm YM}[\xi],
\qquad
\xi^i=\alpha_E\alpha_B\chi^{ij}w_j,}
\label{eq:YM_bracket}
\end{equation}
with
\begin{equation}
    V_{\rm YM}[\xi]
    \equiv\int d^3x\,\xi^i\kappa_{ab}F^a_{ij}E^{jb}.
\end{equation}
This is the gauge-covariant spatial-deformation generator. With the convention
\begin{equation}
    G[\Lambda]
    \equiv-\int d^3x\,\Lambda^a\mathcal D_iE_a^i,
\end{equation}
the generator of the ordinary Lie derivative of the connection is
\begin{equation}
    D_{\rm YM}[\xi]
    =V_{\rm YM}[\xi]+G[\xi^iA_i],
\end{equation}
so Eq.~\eqref{eq:YM_bracket} is equivalently
\begin{equation}
\boxed{
\{H_{\rm YM}[N],H_{\rm YM}[M]\}
=D_{\rm YM}[\xi]-G[\xi^iA_i].}
\label{eq:YM_diff_plus_Gauss}
\end{equation}
The Gauss term is an internal gauge transformation, not a second spacetime structure function. Canonical Maxwell theory exhibits the same distinction between hypersurface deformation and the preserved Gauss constraint \cite{KucharStone1987}.

Progression closure requires the tangential structure function to be the same $h^{ij}$ already selected by the gravitational sector. Since $N$ and $M$ are arbitrary,
\begin{equation}
\boxed{
\alpha_E\alpha_B\chi^{ij}=h^{ij}.}
\label{eq:YM_closure_condition}
\end{equation}
Thus the physical gauge-field co-metric is $h^{ij}$. If the trial tensor is normalized as $\chi^{ij}=h^{ij}$, closure requires
\begin{equation}
    \alpha_E\alpha_B=1.
\end{equation}
The ordinary Yang--Mills normalization
\begin{equation}
    \alpha_E=g^2,
    \qquad
    \alpha_B=g^{-2}
\end{equation}
automatically satisfies this condition. The internal gauge coupling $g$ is therefore not fixed by progression closure; only the reciprocal electric/magnetic normalization required for a common causal cone is fixed.

For the closed branch $\chi^{ij}=h^{ij}$, the Yang--Mills Hamiltonian contains no gravitational momentum $\pi^{ij}$ and depends algebraically on $h_{ij}$. Its antisymmetrized gravity--gauge cross bracket is consequently proportional to $NM-MN$ and vanishes. The gravity plus gauge sector therefore closes as
\begin{equation}
\boxed{
\{H_{g+{\rm YM}}[N],H_{g+{\rm YM}}[M]\}
=D_g[h^{ij}w_j]+V_{\rm YM}[h^{ij}w_j],}
\label{eq:gravity_YM_closure}
\end{equation}
or, equivalently, as the ordinary total spatial diffeomorphism plus the corresponding Gauss transformation.

\subsection{Dirac matter: principal/characteristic matching and the local Lorentz constraint}
\label{subsec:Dirac_closure}

Fermions require a separate treatment because the Dirac action is first order in progression derivatives. The canonical spinor momenta therefore generate second-class constraints. The appropriate reduced bracket is a Dirac bracket, and a tetrad or spatial-triad lift of $h_{ij}$ introduces local Lorentz gauge freedom. This is standard in canonical gravity coupled to electrons, where the Hamiltonian constraints close using Dirac brackets together with local Lorentz generators \cite{NelsonTeitelboim1978}. The triad is gauge data for the same metric $h_{ij}$ rather than a second physical spatial geometry.

After solving the fermionic second-class constraints, use a spatial half-density spinor $\Psi$ with reduced bracket, up to the conventional Grassmann sign and factor,
\begin{equation}
    \{\Psi_A(x),\Psi_B^\dagger(y)\}_D
    =-i\delta_{AB}\delta^3(x-y).
\end{equation}
To test the characteristic geometry without assuming it, introduce matrices
\begin{equation}
    \alpha^i=r_a{}^i\alpha^a,
    \qquad
    \{\alpha^i,\alpha^j\}=2\chi_D^{ij},
    \qquad
    \chi_D^{ij}\equiv\delta^{ab}r_a{}^ir_b{}^j,
\end{equation}
where $\{\alpha^a,\alpha^b\}=2\delta^{ab}$. Allow a dimensionless propagation coefficient $v$. A Hermitian normal generator is
\begin{equation}
\boxed{
H_D[N]=\int d^3x\,N\left[
-\frac{iv}{2}\left(
\Psi^\dagger\alpha^i\mathcal D_i\Psi
-(\mathcal D_i\Psi)^\dagger\alpha^i\Psi
\right)
+m\Psi^\dagger\beta_D\Psi
\right].}
\label{eq:Dirac_trial_H}
\end{equation}
Here $\mathcal D_i$ contains the spin connection and any internal gauge connection. After integrating by parts, the one-particle normal operator can be written
\begin{equation}
    K_N=-ivP_N+Nm\beta_D,
    \qquad
    P_N=N\alpha^i\mathcal D_i+\frac12\alpha^iD_iN.
    \label{eq:Dirac_PN}
\end{equation}
The $\frac12\alpha^iD_iN$ term is essential. At a point choose a local frame for the principal-symbol calculation so that $\mathcal D_i\alpha^j=0$. Expanding $[P_N,P_M]$, the apparent second-derivative terms cancel and the first-derivative coefficient is
\begin{equation}
\begin{aligned}
[P_N,P_M]_{\rm 1st}
&=w_i\alpha^i\alpha^j\mathcal D_j
-\frac12w_i[\alpha^i,\alpha^j]\mathcal D_j\\
&=\frac12w_i\{\alpha^i,\alpha^j\}\mathcal D_j
=\chi_D^{ij}w_i\mathcal D_j.
\end{aligned}
\label{eq:Dirac_principal_bracket}
\end{equation}
The mass does not alter this structure function. Using $\{\alpha^i,\beta_D\}=0$, the cross terms satisfy
\begin{equation}
    [P_N,M\beta_D]+[N\beta_D,P_M]=0,
    \label{eq:Dirac_mass_cancel}
\end{equation}
and $[Nm\beta_D,Mm\beta_D]=0$. Thus the mass is unrestricted by the deformation algebra.

Restoring the spin and gauge connections adds zero-order internal rotations, but those lower-order terms have not been recomputed here for an initially independent trial $\chi_D^{ij}$. What the explicit calculation establishes is the characteristic tangential coefficient
\begin{equation}
    \boxed{\xi^i_{\rm pr}=v^2\chi_D^{ij}w_j.}
    \label{eq:Dirac_structure_function}
\end{equation}
The standard gravity--electron analysis of Nelson and Teitelboim shows that, once the ordinary tetrad geometry is adopted, the completed constrained system closes with spatial, local-Lorentz, and internal-gauge generators \cite{NelsonTeitelboim1978}. That known completion strongly supports the $\chi_D^{ij}=h^{ij}$ branch, but it does not replace a full lower-order Dirac-bracket calculation for an arbitrary independent trial $\chi_D^{ij}$.

Matching the \emph{principal/characteristic} progression-HDA structure gives
\begin{equation}
\boxed{
    v^2\chi_D^{ij}=h^{ij}.}
\label{eq:Dirac_closure_condition}
\end{equation}
If $\chi_D^{ij}=h^{ij}$ is used as the normalized candidate co-metric, the forward-progression branch gives $v=1$. The explicit result is therefore that Dirac matter shares the same \emph{principal characteristic cone} as the scalar and gauge sectors, while $m$ remains free. Full closure of the independent trial Dirac sector remains a separate lower-order constraint calculation.

For charged spinors the internal gauge constraint contains both the electric divergence and the fermion charge density,
\begin{equation}
    \mathcal G_a
    =-\mathcal D_iE_a^i+\rho_a^{\rm f}\approx0.
\end{equation}
The gauge terms generated separately by the Yang--Mills and Dirac brackets therefore combine into this same total Gauss constraint rather than introducing another spacetime deformation rule.

\subsection{Combined matter cones and common lapse smearing}
\label{subsec:combined_matter_closure}

The scalar, gauge, and fermionic calculations may be summarized by the three physical structure-function conditions
\begin{equation}
\boxed{
\mathsf A^{CA}\mathsf C^{ij}_{AB}=\delta^C{}_Bh^{ij},
\qquad
\alpha_E\alpha_B\chi^{ij}=h^{ij},
\qquad
v^2\chi_D^{ij}=h^{ij}.}
\label{eq:three_matter_closures}
\end{equation}
Thus the principal Standard-Model field types tested here share a single physical spatial co-metric and characteristic cone. Their internal parameters remain free: scalar potentials and field-space metrics, Yang--Mills couplings, and fermion masses do not enter the common HDA structure function.

For gravity plus scalar and Yang--Mills fields, the brackets calculated here close completely with the single spacetime structure function $h^{ij}$, modulo the Gauss constraint. For the Dirac sector the present calculation establishes the same structure function at principal/characteristic order. On the standard tetrad--spinor branch, the known completed gravity--electron theory has the schematic first-class form
\begin{equation}
\boxed{
\{H_{\rm tot}[N],H_{\rm tot}[M]\}_D
=D_{\rm tot}[h^{ij}w_j]
+G_{\rm tot}[\Lambda_{NM}]
+J_{\rm Lor}[\Theta_{NM}],}
\label{eq:full_matter_HDA}
\end{equation}
but Eq.~\eqref{eq:full_matter_HDA} should not be read as an independent full-closure proof for an arbitrary trial $\chi_D^{ij}$.

The normal deformation generator is linear in the arbitrary canonical smearing $N$. If, only after the canonical algebra has been established, one chooses an adapted operational representation in which a physical progression profile is represented by the lapse smearing, then a replacement $N\mapsto N^q$ would be incompatible with that same linear generator unless $q=1$ (or unless the generator itself were changed). This is only a \emph{consistency check} on the earlier one-factor temporal rate/measure argument; it is not an independent derivation of a pregeometric weighting exponent, because the identification between a progression representative and a canonical smearing is made downstream.

A species-dependent rescaling of the lapse coupling can be tested only after its normalization has been fixed independently. Formally, smearing species $I$ by $\zeta_I N$ makes its self-bracket carry $\zeta_I^2 h^{ij}$, but constant factors may in some matter parametrizations be absorbed into field or coupling normalizations. We therefore do not interpret $\zeta_I^2=1$ as an independent universality theorem here. The robust closure result is the common principal characteristic co-metric; stronger claims about species-dependent progression charges require an operational definition that fixes their normalization.

\subsection{Reconstructed scalar action and covariant interpretation}

Define the scalar field-space metric
\begin{equation}
    G_{AB}\equiv(\mathsf A^{-1})_{AB}.
\end{equation}
Using Eq.~\eqref{eq:C_metric_factorization}, the canonical scalar action is
\begin{equation}
    S_m=\int d\lambda\,d^3x
    \left[p_A\dot\phi^A-N\mathcal H_m-N^ip_AD_i\phi^A\right].
\end{equation}
Eliminating $p_A$ gives
\begin{equation}
\begin{aligned}
S_m=\int d\lambda\,d^3x\,N\sqrt h\Bigg[&
\frac{1}{2N^2}G_{AB}
(\dot\phi^A-N^iD_i\phi^A)
(\dot\phi^B-N^jD_j\phi^B)\\
&-\frac12G_{AB}h^{ij}D_i\phi^A D_j\phi^B-U(\phi)
\Bigg].
\end{aligned}
\label{eq:matter_ADM_action}
\end{equation}
With
\begin{equation}
    ds^2=-N^2d\lambda^2
    +h_{ij}(dx^i+N^id\lambda)(dx^j+N^jd\lambda),
    \qquad \sqrt{-g}=N\sqrt h,
\end{equation}
this is the ordinary minimally metric-coupled scalar action
\begin{equation}
\boxed{
S_m=-\int d^4x\sqrt{-g}
\left[
\frac12G_{AB}(\phi)g^{\mu\nu}
\partial_\mu\phi^A\partial_\nu\phi^B+U(\phi)
\right],}
\label{eq:reconstructed_covariant_matter}
\end{equation}
up to the sign convention associated with the metric signature and scalar kinetic term. The gauge closure condition and the calculated Dirac principal/characteristic matching select the same $N$ and $h_{ij}$ at the respective levels established here. On the standard tetrad--spinor branch, the additional Gauss and local Lorentz constraints express internal redundancy rather than another spacetime geometry. The covariant stress--energy identity $\nabla_\mu T^{\mu\nu}=0$ then follows on the matter equations of motion for the complete minimally coupled system.

The matter closure established here is still conditional on the chosen local field classes. Nonminimal curvature couplings, higher derivatives, composite effective matter, and enlarged physical phase spaces require separate analysis. In particular, the result obtained for the independent trial Dirac sector is only principal/characteristic matching. The standard tetrad--spinor constrained theory supplies a known lower-order completion, but this paper does not rederive full closure for arbitrary trial $\chi_D^{ij}$.

\section{Stress--energy conservation after matter closure}

The scalar and gauge closure calculations together with the Dirac principal/characteristic calculation in Secs.~\ref{subsec:matter_closure}--\ref{subsec:combined_matter_closure} supply a concrete route from progression HDA closure to one common matter geometry at the levels established here. The scalar action is reconstructed explicitly, Maxwell/Yang--Mills closes modulo its Gauss constraint, and the trial Dirac sector is shown to share the same principal co-metric; the standard tetrad--spinor branch supplies the known lower-order Lorentz/gauge completion. Under the corresponding minimal covariant completion this gives $\nabla_\mu T^{\mu\nu}=0$ for the total matter system. The present section records a complementary observer-current interpretation and clarifies why a single temporal rate/measure factor, by itself, is insufficient to imply the four local tensor conservation equations.

\subsection{Observer energy currents}

For a local observer with unit timelike four-velocity $u^\nu$, define the measured energy current
\begin{equation}
    \boxed{J^\mu[u]=-T^{\mu\nu}u_\nu.}
    \label{eq:observer_current}
\end{equation}
The relevant local premise is not that $J^\mu[u]$ is globally divergence-free for an arbitrary accelerating observer. Even when $\nabla_\mu T^{\mu\nu}=0$, one has
\begin{equation}
    \nabla_\mu J^\mu[u]
    =-T^{\mu\nu}\nabla_\mu u_\nu
\end{equation}
for a general congruence. The useful statement is pointwise: at any event $p$, choose a freely falling frame and extend $u^\nu$ locally so that
\begin{equation}
    \nabla_\mu u_\nu\big|_p=0.
\end{equation}
If the progression-referenced local energy balance holds for every such timelike $u^\nu$, then
\begin{equation}
    0=\nabla_\mu J^\mu[u]\big|_p
    =-u_\nu\nabla_\mu T^{\mu\nu}\big|_p.
\end{equation}
Because this must hold for every timelike direction at $p$,
\begin{equation}
    \boxed{\nabla_\mu T^{\mu\nu}=0.}
    \label{eq:T_conservation}
\end{equation}
The logical strengthening is therefore
\begin{equation}
    \boxed{
    \text{one-factor temporal rate/measure conversion}
    \ +\ 
    \text{local energy balance for every inertial observer}.}
    \label{eq:strengthening}
\end{equation}
Under this stronger universal premise, tensor conservation follows pointwise once the invariant measure and genuine spacetime currents have emerged.

\subsection{Matter-action formulation}

There is a complementary and more systematic route. If Eq.~\eqref{eq:measure_step} extends to a local matter action
\begin{equation}
    \boxed{
    S_m[g,\psi]
    =\int d^4x\sqrt{-g}\,\mathcal L_m(g,\psi,\nabla\psi),}
    \label{eq:matter_action}
\end{equation}
with no additional preferred temporal structure, define
\begin{equation}
    T_{\mu\nu}
    =-\frac{2}{\sqrt{-g}}
      \frac{\delta S_m}{\delta g^{\mu\nu}}.
\end{equation}
The diffeomorphism Noether identity then gives Eq.~\eqref{eq:T_conservation} on the matter equations of motion. In this formulation the key progression-specific task is to derive why local matter physics depends only on the emergent invariant structures rather than separately on $\lambda$ or an external preferred frame.

The two arguments are consistent. The observer-current construction exhibits the physical content of the conservation principle, while the action formulation supplies the compact covariant implementation.

\section{Relation to self-coupling and equilibrium closure routes}

The canonical calculation above is the operative reconstruction in this paper. Two established lines of work are nevertheless useful for interpreting, rather than deriving, the resulting nonlinear completion.

First, perturbative self-coupling arguments show why a universally coupled geometric response cannot consistently remain an external field sourced only by matter. If a linear response obeys $\mathcal E^{(1)}_{\mu\nu}=\kappa T^m_{\mu\nu}$, nonlinear geometric terms may be moved to the source side as a background-dependent $t^{\rm resp}_{\mu\nu}$ so that the conserved perturbative source is $T^m_{\mu\nu}+t^{\rm resp}_{\mu\nu}$. Iterating universal massless spin-2 self-coupling closes on Einstein gravity under the usual locality and gauge-consistency assumptions \cite{Deser1970,Deser2010}. This supports the progression interpretation that the response participates in the same total balance, while not implying a unique generally covariant local gravitational stress tensor.

Second, Jacobson-type arguments derive Einstein dynamics from local equilibrium conditions rather than from a canonical bootstrap. The Clausius relation on local acceleration horizons, together with Unruh temperature, area entropy, Raychaudhuri focusing, and stress--energy conservation, yields the Einstein equation as an equation of state \cite{Jacobson1995}; causal-diamond and entanglement-equilibrium formulations provide related local balance principles \cite{Jacobson2016,JacobsonVisser2019}. These results suggest a possible future progression-adapted equilibrium construction, but the required entropy functional and its relation to the phase-defined $d\sigma$ are not derived here. Likewise, the non-equilibrium entropy-production structures discussed by Eling, Guedens, and Jacobson provide only an analogy for the quadratic static progression stiffness derived below \cite{Eling2006}.

These routes therefore serve as comparison frameworks. The present paper does not use them to establish the HDA or the canonical dynamics; its substantive upstream claim remains the conditional progression-integrability route developed in Sec.~\ref{sec:path_independence}.

\section{Static two-function reduction of the canonical completion}

The canonical progression-HDA calculation selects the ADM constraint and reconstructs the corresponding geometric action within the minimal phase-space assumptions. We now reduce that result in the static isotropic sector. This section therefore functions as a consistency check and phenomenological bridge rather than as an independent Einstein--Hilbert assumption.

\subsection{Metric ansatz}

Take
\begin{equation}
    \boxed{
    ds^2=-e^{2A(\mathbf x)}c^2d\lambda^2
    +e^{2B(\mathbf x)}\delta_{ij}dx^idx^j,}
    \label{eq:AB_metric}
\end{equation}
with
\begin{equation}
    A=\ln\Pi.
\end{equation}
Here $A$ is the progression/lapse response and $B$ is an independent isotropic spatial response. In ADM language
\begin{equation}
    N=e^A,\qquad N^i=0,\qquad h_{ij}=e^{2B}\delta_{ij}.
\end{equation}
The configuration is static, so $K_{ij}=0$.

For the conformally flat spatial metric,
\begin{equation}
    {}^{(3)}R
    =e^{-2B}\left[-4\nabla^2B-2|\nabla B|^2\right].
    \label{eq:R3}
\end{equation}
Since $N\sqrt h=e^{A+3B}$, integration by parts gives the exact static bulk action
\begin{equation}
    \boxed{
    S_g=\frac{c^4}{16\pi G}
    \int d\lambda\,d^3x\,
    e^{A+B}
    \left[4\nabla A\cdot\nabla B+2|\nabla B|^2\right].}
    \label{eq:AB_action}
\end{equation}
Boundary terms are assumed to be treated in the standard way.

\subsection{The variable \texorpdfstring{$S=A+B$}{S=A+B}}

Define
\begin{equation}
    S\equiv A+B.
\end{equation}
Then
\begin{equation}
    4\nabla A\cdot\nabla B+2|\nabla B|^2
    =2|\nabla S|^2-2|\nabla A|^2,
\end{equation}
so
\begin{equation}
    \boxed{
    S_g=\frac{c^4}{8\pi G}
    \int d\lambda\,d^3x\,
    e^S\left(|\nabla S|^2-|\nabla A|^2\right).}
    \label{eq:AS_action}
\end{equation}
The earlier one-function ansatz $B=-A$ corresponds to freezing $S=0$. The two-function reduction allows $S$ to be generated dynamically by nonlinear progression gradients.

\subsection{Matter coupling and coupled static equations}

For a prescribed static nonrelativistic source with coordinate rest-mass density $\rho_c$,
\begin{equation}
    S_m=-\int d\lambda\,d^3x\,\rho_c c^2 e^A.
    \label{eq:dust_action}
\end{equation}
Variation of $A$ and $S$ gives
\begin{equation}
    \boxed{
    \nabla\cdot(e^S\nabla A)
    =\frac{4\pi G}{c^2}\rho_c e^A,}
    \label{eq:A_eq}
\end{equation}
\begin{equation}
    \boxed{
    2\nabla^2S+|\nabla S|^2+|\nabla A|^2=0}
    \label{eq:S_eq}
\end{equation}
for this prescribed static nonrelativistic source. With relativistic pressure and stress, the full matter action contributes to the spatial equations as expected.

Equation~\eqref{eq:S_eq} displays the self-response directly: even where the matter density vanishes, a nonzero progression gradient can source the nonlinear spatial relation between $A$ and $B$.

\subsection{Weak-field limit}

Expand
\begin{equation}
    A=A_1+A_2+\cdots,
    \qquad
    S=S_1+S_2+\cdots.
\end{equation}
At first order, Eq.~\eqref{eq:S_eq} gives
\begin{equation}
    \nabla^2S_1=0.
\end{equation}
For an isolated system with $S_1\to0$ at infinity,
\begin{equation}
    \boxed{S_1=0,\qquad B_1=-A_1.}
    \label{eq:BminusA}
\end{equation}
Equation~\eqref{eq:A_eq} then becomes
\begin{equation}
    \boxed{
    \nabla^2A_1=\frac{4\pi G}{c^2}\rho.}
    \label{eq:PoissonA}
\end{equation}
Identifying
\begin{equation}
    A_1=\frac{\Phi}{c^2}
\end{equation}
recovers
\begin{equation}
    \nabla^2\Phi=4\pi G\rho.
\end{equation}
For a positive isolated mass, $\Phi=-GM/r$, so $A<0$ and
\begin{equation}
    \Pi=e^A<1.
\end{equation}
Thus positive gravitating energy corresponds to reduced local progression relative to infinity.

At quadratic order, Eq.~\eqref{eq:AS_action} with $S_1=0$ yields the effective progression action
\begin{equation}
    \boxed{
    S_{\rm eff}^{(2)}[A]
    =-\frac{c^4}{8\pi G}
    \int d\lambda\,d^3x\,|\nabla A|^2.}
    \label{eq:A_eff_action}
\end{equation}
The corresponding static functional contains
\begin{equation}
    \boxed{
    E_{\rm grad}^{(2)}
    =\frac{c^4}{8\pi G}\int d^3x\,|\nabla A|^2.}
    \label{eq:grad_energy}
\end{equation}
Since $A=\ln\Pi=\delta\Pi+O(\delta\Pi^2)$, this is precisely the weak-field quadratic progression stiffness previously required to generate a Poisson profile from a localized energetic excess.

At second order in vacuum,
\begin{equation}
    \boxed{
    \nabla^2S_2=-\frac12|\nabla A_1|^2.}
    \label{eq:S2}
\end{equation}
Thus the progression gradient itself generates the nonlinear spatial response. For $A_1=-GM/(rc^2)\equiv-U$,
\begin{equation}
    S_2=-\frac14U^2,
    \qquad
    B=U-\frac14U^2+O(U^3).
\end{equation}

\subsection{Exact spherical vacuum solution}

The two-function reduction can be solved exactly in static spherical vacuum. The lapse variation is equivalent to
\begin{equation}
    2\nabla^2B+|\nabla B|^2=0.
\end{equation}
With
\begin{equation}
    \psi=e^{B/2},
\end{equation}
this becomes
\begin{equation}
    \nabla^2\psi=0.
\end{equation}
Asymptotic flatness and spherical symmetry give
\begin{equation}
    \psi=1+\frac{GM}{2rc^2}.
\end{equation}
The remaining vacuum equation integrates to
\begin{equation}
    \boxed{
    e^A=\Pi
    =\frac{1-\dfrac{GM}{2rc^2}}
           {1+\dfrac{GM}{2rc^2}},}
    \label{eq:Pi_exact}
\end{equation}
\begin{equation}
    \boxed{
    e^B=\left(1+\frac{GM}{2rc^2}\right)^2.}
    \label{eq:B_exact}
\end{equation}
Hence
\begin{equation}
    ds^2=-\left(
    \frac{1-GM/(2rc^2)}{1+GM/(2rc^2)}
    \right)^2c^2d\lambda^2
    +\left(1+\frac{GM}{2rc^2}\right)^4d\mathbf x^2,
    \label{eq:schwarzschild_iso}
\end{equation}
which is Schwarzschild spacetime in isotropic coordinates. The one-function approximation $B=-A$ is therefore correct only at leading weak-field order; the independent spatial response is essential nonlinearly.

\subsection{Maxwell propagation and the former \texorpdfstring{$p=2$}{p=2} optical closure}

For a null ray in Eq.~\eqref{eq:AB_metric},
\begin{equation}
    c\,d\lambda=e^{B-A}dl,
\end{equation}
so the isotropic coordinate optical index is
\begin{equation}
    \boxed{n_{\rm eff}=e^{B-A}.}
    \label{eq:neff}
\end{equation}
At first order $B=-A$, giving
\begin{equation}
    \boxed{
    n_{\rm eff}=e^{-2A}+O(A^2)
    =\Pi^{-2}+O(A^2).}
    \label{eq:p2}
\end{equation}
The Maxwell/Yang--Mills closure result of Sec.~\ref{subsec:YM_closure} makes this more than a metric dictionary: the gauge-field Hamiltonian itself is forced to use the same physical co-metric $h^{ij}$. The weak-field optical response is therefore not an independently tuned constitutive index. One factor arises from the lapse $N=e^A=\Pi$ and the second from the HDA-related spatial response $B=-A$ at leading order. The relation $n_{\rm eff}\simeq\Pi^{-2}$ is consequently a weak-field consequence of universal Maxwell closure plus the static spatial response, not an exact nonlinear law. For the exact spherical solution,
\begin{equation}
    n_{\rm eff}
    =\frac{(1+U/2)^3}{1-U/2}
    =1+2U+\frac74U^2+O(U^3).
\end{equation}
Massless propagation therefore belongs to the full temporal-spatial geometry rather than to an $m\to0$ limit of the massive rest-energy phase.

\section{Quantum phase as the operational matter response}

The geometric completion determines proper time. Quantum phase then gives an operational microscopic description of how matter responds to that proper-time structure. This response should not be added as an independent force on top of geodesic motion; it is the wave-mechanical manifestation of the same coupling.

\subsection{Weak-field matter Hamiltonian}

For a massive particle in the static metric \eqref{eq:AB_metric},
\begin{equation}
    S_p=-mc^2\int d\tau.
\end{equation}
For slow coordinate velocity,
\begin{equation}
    d\tau
    =e^A d\lambda
    \sqrt{1-e^{2(B-A)}\frac{v^2}{c^2}},
\end{equation}
so
\begin{equation}
    L
    =-mc^2e^A
    +\frac{m}{2}e^{2B-A}v^2+O(v^4/c^2).
    \label{eq:particle_L}
\end{equation}
To first weak-field order this is
\begin{equation}
    L=-mc^2+\frac12mv^2-mc^2A+O(A v^2,A^2).
\end{equation}
Thus the leading Hamiltonian is
\begin{equation}
    \boxed{
    H_\lambda
    =\frac{p^2}{2m}+mc^2A
    =\frac{p^2}{2m}+mc^2\ln\Pi}
    \label{eq:H_A}
\end{equation}
up to the irrelevant rest-energy constant and higher-order metric corrections.

This refines the earlier heuristic use of $mc^2(\Pi-1)$. Since
\begin{equation}
    \Pi-1=A+O(A^2),
\end{equation}
both forms agree in the Newtonian limit, but $A=\ln\Pi$ is the natural nonlinear variable once the geometric completion is taken seriously.

\subsection{Ehrenfest acceleration}

The weak-field Schr\"odinger equation is
\begin{equation}
    i\hbar\frac{\partial\psi}{\partial\lambda}
    =\left[-\frac{\hbar^2}{2m}\nabla^2+mc^2A\right]\psi.
\end{equation}
Ehrenfest's theorem gives
\begin{equation}
    \frac{d\langle\mathbf p\rangle}{d\lambda}
    =-mc^2\langle\nabla A\rangle.
\end{equation}
For a narrow nonrelativistic packet,
\begin{equation}
    \boxed{
    \mathbf a
    =-c^2\nabla A
    =-c^2\nabla\ln\Pi.}
    \label{eq:a_lnPi}
\end{equation}
The test mass cancels. With $A=\Phi/c^2$,
\begin{equation}
    \mathbf a=-\nabla\Phi.
\end{equation}
The phase picture therefore reproduces the same leading motion as the geometric particle action while retaining the standard Born rule.

\subsection{Clock phase and motion are one coupling, not two}

The local clock relation
\begin{equation}
    d(\Delta\theta)=-\frac{\Delta E}{\hbar}d\tau
\end{equation}
and the particle action
\begin{equation}
    S_p=-mc^2\int d\tau
\end{equation}
are both consequences of the same proper-time coupling. This is important conceptually. The framework need not postulate one law for clock redshift and a second law for free fall. If all local matter Hamiltonians evolve with the same $d\tau$, internal energy, rest energy, and composite binding energy participate through a common geometric/progression structure. Experimental universality then becomes a direct constraint on any proposed departure from that common coupling.

\section{Conservation, equilibrium, and progression redistribution}

The original motivating intuition was that an inhomogeneous energetic environment may require a spatial redistribution of progression in order to preserve a deeper conservation relation. The present development sharpens that intuition by distinguishing three levels.

First, the kinematic conversion
\begin{equation}
    dE/d\lambda=\Pi\,dE/d\tau
\end{equation}
selects the $q=1$ measure factor. Second, local conservation of energy transfer for all observers leads, after the measure is geometrized, to $\nabla_\mu T^{\mu\nu}=0$. Third, the actual spatial profile of progression requires either geometric dynamics or an equilibrium principle. Conservation alone does not specify a $1/r$ profile.

The weak-field variational functional
\begin{equation}
    E[A]
    =\int d^3x\left[
    \frac{c^4}{8\pi G}|\nabla A|^2
    +\rho c^2 A\right]
    \label{eq:weak_functional}
\end{equation}
can now be interpreted as the static quadratic reduction of the two-function geometric completion rather than as a canonical scalar-field energy postulated from the start. Its variation yields Eq.~\eqref{eq:PoissonA}. The foundational question has therefore shifted from ``why choose a gradient term?'' to ``why should the progression-weighted conservation principle close on the geometric dynamics whose weak-field reduction contains that term?''

A stronger equilibration statement such as
\begin{equation}
    \Pi E_{\rm loc}\rightarrow\text{constant}
\end{equation}
may still be meaningful for particular connected systems, but it is not required for the local gravitational derivation and should not be identified with the universal conservation law.

\section{Relation to established quantum and gravitational results}

Several parts of the framework are standard physics expressed in a different causal order. Matter-wave phase depends on proper time and action \cite{Stodolsky1979}; relativistic atom interferometry systematically computes gravitational phase shifts \cite{Dimopoulos2008}; quantum-clock interferometry isolates proper-time differences through internal-state phase \cite{Zych2011,Roura2020,Roura2021}; and joint tests of gravitational redshift and free fall directly constrain nonuniversal clock/motion couplings \cite{Ufrecht2020}.

The integrability principle used here has direct precedents in relativistic quantum theory. Tomonaga and Schwinger formulated quantum evolution as a local functional evolution of states on spacelike hypersurfaces, with consistency requiring independence from arbitrary compatible local deformations \cite{Tomonaga1946,Schwinger1948}. In ordinary QFT this is tied to microcausality at spacelike separation. Schwinger's local energy--momentum-density commutator gives a sharper canonical analogue: derivative-of-delta momentum-density terms survive at coincidence, and after smearing generate the antisymmetric combination $N\partial_iM-M\partial_iN$ characteristic of normal-deformation commutators \cite{Schwinger1963}. Fukuyama explicitly derived the energy-density commutation relation from an integrability condition for local Hamiltonian evolution \cite{Fukuyama1973}. These results presuppose spacetime/hypersurface structure, so they motivate but do not derive the pregeometric postulate of Sec.~\ref{sec:path_independence}.

The bridge to canonical gravity is especially close. Moncrief and Teitelboim showed that the momentum constraints of general relativity follow as integrability conditions when the Hamiltonian constraint is imposed on the Hamilton principal functional or quantum state functional, so that the content of the theory may be expressed through a Tomonaga--Schwinger-type functional equation \cite{MoncriefTeitelboim1972}. Teitelboim's subsequent path-independence analysis and the Hojman--Kucha\v{r}--Teitelboim reconstruction are established precedents for the deformation-algebra route \cite{Teitelboim1973,HKT1976}. The present progression construction is intended to begin one logical step earlier: an operational clock notion motivates local relational progression, and the Tomonaga--Schwinger tradition motivates integrability of that local evolution; the spatial structure function is then derived from the pregeometric locality assumptions rather than presupposed from an embedding. The subsequent canonical uniqueness calculation is not claimed as a new general proof of $\mathrm{HDA}\Rightarrow\mathrm{GR}$; it is an explicit restricted reconstruction within the stated minimal class.

A recent relational quantum construction provides a complementary precedent. H\"ohn, Russo, and Smith formulate matter relative to quantum hypersurfaces and show that the conditional matter wave functional satisfies the Tomonaga--Schwinger equation \cite{Hoehn2024}. This supports the logical compatibility of relational quantum descriptions with many-fingered local evolution, but it does not establish the present pregeometric progression algebra.

Another nearby relational program is the ``3-space'' or ``relativity without relativity'' approach of Barbour, Foster, and N. {\'O} Murchadha, which seeks GR without presupposing spacetime or the relativity principle and uses constraint propagation to select consistent geometrodynamics and common matter propagation speeds \cite{Barbour2002}. Anderson and Barbour extended that mechanism to interacting vector fields and Yang--Mills structure \cite{AndersonBarbour2002}. The distinction is methodological: the 3-space program begins from a restricted relational action on configuration space and lets constraint propagation select viable theories; the present construction begins from operational local progression and integrability modulo spatial labels, derives the candidate deformation structure, and only then asks which canonical gravity and matter Hamiltonians realize it.

A closely related modern program is the ``gravitational closure'' approach of D\"ull, Schuller, Stritzelberger, and Wolz \cite{Duell2018}. There the causal/principal structure of specified matter equations is taken as input, and foliation-independent joint canonical evolution is used to derive compatible gravitational dynamics; for Standard-Model matter the resulting closure gives general relativity. The direction of explanation proposed here is different: relational progression integrability is used to motivate a spatial structure function and common HDA first, after which scalar, gauge, and fermionic sectors are tested against that geometry. The two programs are therefore complementary rather than competing, and the gravitational-closure literature provides an important benchmark for the matter-cone claims made here.

The gauge bracket closes modulo its Gauss constraint, consistent with canonical Maxwell analyses of hypersurface deformations \cite{KucharStone1987}; the fermion calculation establishes principal/characteristic matching and uses the standard Dirac-bracket/tetrad framework in which local Lorentz constraints accompany spacetime covariance \cite{NelsonTeitelboim1978}. Modern canonical discussions also emphasize that the hypersurface-deformation algebra and spacetime diffeomorphisms have a subtle off-shell relation, especially when deformation parameters become phase-space dependent \cite{Thiemann2024,Bojowald2025}. For that reason this paper keeps the arbitrary canonical smearings $N,N^i$ distinct from operational progression observables until an adapted representation is chosen. The spin-2 self-coupling and divergence-free geometric routes remain valuable complementary interpretations of nonlinear self-consistency \cite{Deser1970,Deser2010,Lovelock1971}. Jacobson's horizon thermodynamics, entanglement equilibrium, and causal-diamond work provide an alternative energy-balance route to the Einstein equation \cite{Jacobson1995,Jacobson2016,JacobsonVisser2019}.

\section{A staged derivation program}

The calculations now separate what has been obtained conditionally from what still must be derived from the underlying progression premise.

\begin{enumerate}[label=\textbf{Step \arabic*:},leftmargin=1.05in]
    \item \textbf{Progression covariance.} Establish microscopically that $\lambda$ is arbitrary and that observables depend only on the phase-defined clock increment $d\sigma$, relational clock ratios, and ultimately their downstream proper-time representation.

    \item \textbf{Local-evolution hypothesis.} Derive why independently specifiable local progression defines a many-fingered relational evolution rather than merely a set of clock comparisons. A quantum implementation should identify the appropriate local generators and physical state space without presupposing the spacetime geometry to be derived.

    \item \textbf{Relational progression integrability and leading locality.} Derive, rather than assume, the state-space bundle structure and the vertical-curvature condition of Eq.~\eqref{eq:vertical_curvature}. Tomonaga--Schwinger and Dirac--Schwinger integrability provide precedents, not a pregeometric proof. Determine independently whether the leading local commutator is single-gradient and whether normal and tangential generators are independent.

    \item \textbf{Kinematic progression lemma.} Under those assumptions, derive Eq.~\eqref{eq:kinematic_lemma}; use Jacobi closure to obtain symmetry of $\beta^{ij}$ and Eq.~\eqref{eq:jacobi_allowed_beta_response}.

    \item \textbf{Spatial co-metric and pointwise locality.} Derive positivity/nondegeneracy and determine whether pointwise progression locality removes the Jacobi-allowed $C^{(ijk)}\partial_kN$ term.

    \item \textbf{Canonical gravitational closure.} Within the minimal $(h_{ij},\pi^{ij})$ phase space, the calculation of Sec.~\ref{sec:canonical_closure} gives $b=-a/2$, $ac=-1$, and Eq.~\eqref{eq:K_pi_relation}. The remaining foundational question is why this minimal canonical class follows from progression conservation.

    \item \textbf{Potential uniqueness.} The four-spatial-derivative test eliminates $R^2$, $R_{ij}R^{ij}$, and $D^2R$ within the undeformed HDA. Extend this analysis to parity-odd terms, enlarged phase spaces, and any progression-motivated higher-derivative structures before claiming uniqueness beyond the tested class.

    \item \textbf{Matter closure.} The scalar and Maxwell/Yang--Mills closure calculations and the Dirac principal/characteristic calculation give Eq.~\eqref{eq:three_matter_closures}. Gauge closure is modulo the Gauss constraint; the fermion calculation uses the reduced Dirac bracket only to establish principal/characteristic matching, while the lower-order Lorentz/gauge completion is known for the standard tetrad--spinor branch. The remaining universality tests are composite and bound-state matter, nonminimal curvature couplings, and higher-derivative effective sectors.

    \item \textbf{Measure and tensor conservation.} Derive one power of progression as the temporal rate/measure conversion in Eq.~\eqref{eq:rate_conversion}; separately derive the spatial flux and volume structure needed for a genuine local current balance. Once an adapted lapse representative and $h_{ij}$ emerge, show that the temporal factor belongs in $\sqrt{-g}=N\sqrt h$ and use the completed matter geometry to establish the covariant stress--energy identity without an external preferred structure.

    \item \textbf{Coupling normalization.} Fix the remaining overall constant $a$ from a progression-based source normalization or the measured Newtonian limit, and determine whether $G$ can itself be connected to a microscopic progression scale rather than entered empirically.

    \item \textbf{Total self-consistency and equilibrium.} Show how the nonlinear canonical terms encode the participation of the geometric response in total conservation and compare this result with the Deser bootstrap and Jacobson-style local equilibrium routes.

    \item \textbf{Quantum universality.} Gauge-field and Dirac characteristic propagation now share the completed geometry at classical canonical level. Derive composite-matter evolution, interacting clock phases, bound-state universality, and the semiclassical free-fall limit from the same completed matter geometry.

    \item \textbf{Departure from equilibrium.} Only after the equilibrium theory is fixed should controlled additional degrees of freedom, deformed algebras, or non-equilibrium terms be introduced and tested for distinct phenomenology.
\end{enumerate}

The main shift is that $K_{ij}$ and the minimal nonlinear gravitational action are no longer simply open once the canonical HDA assumptions are adopted. The open foundational burden has moved upward: the framework must derive a genuinely local many-fingered progression dynamics, justify relational progression integrability modulo spatial gauge, and then explain why its leading locality and canonical realization select the undeformed HDA and minimal phase space. Conservation remains essential to the matter and total-balance interpretation, but it need not by itself carry the full kinematic burden of producing the deformation algebra.

\section{Empirical and theoretical failure criteria}

The revised framework faces several immediate tests.

\begin{enumerate}[label=(\roman*)]
    \item \textbf{Preferred-structure failure.} If the theory requires an observable preferred $\lambda$ or local rest frame beyond gauge choice and physical matter state, it conflicts with local Lorentz invariance.

    \item \textbf{Relational-integrability failure.} If independently specified local progression does not define a consistent many-fingered evolution, or if a local update loop that is closed in the physical quotient can leave an observable non-gauge residual normal update (for example an ordering-dependent clock-phase holonomy), the integrability postulate is false and the derived HDA does not follow.

    \item \textbf{Closure-order failure.} If the microscopic update law requires leading two-gradient or higher-gradient commutator terms, the kinematic lemma is not unique and the undeformed HDA need not be the correct algebra.

    \item \textbf{Co-metric failure.} If the Jacobi-symmetric $\beta^{ij}$ is degenerate or not positive on physical spatial directions, it cannot define the spatial co-metric used in the canonical construction.

    \item \textbf{Canonical-class failure.} The gravitational uniqueness result assumes only $(h_{ij},\pi^{ij})$, an ultralocal quadratic momentum term, locality, and the undeformed HDA. If progression conservation requires extra canonical fields, explicit lapse derivatives, or a deformed algebra, Eq.~\eqref{eq:ADM_selected_constraint} is not the unique completion.

    \item \textbf{Higher-derivative failure.} The tested parity-even four-derivative corrections are excluded in the minimal class. If physically required higher-curvature terms survive only by enlarging the phase space or altering the algebra, the framework must state that enlarged structure explicitly rather than retain the minimal uniqueness claim.

    \item \textbf{Matter-universality failure.} The scalar and Maxwell/Yang--Mills sectors realize the same canonical lapse smearing and spatial structure function, while the Dirac sector matches that structure at principal/characteristic order; Gauss/local-Lorentz terms are treated as internal constraints on the corresponding standard branches. The universality claim would fail if composite, bound-state, nonminimal, or higher-derivative matter required a second physical co-metric or a second physical co-metric. Species-dependent normalization factors require an independently fixed operational meaning before they constitute a distinct universality failure.

    \item \textbf{Measure failure.} If the one-factor temporal rate/measure conversion cannot be derived consistently, if the spatial flux/measure needed for a genuine local balance law cannot be obtained, or if a second independent $\Pi$ factor remains after the adapted temporal representative is encoded in $N$ and hence in $\sqrt{-g}$, the conservation interpretation fails or double-counts progression.

    \item \textbf{Optical failure.} Massive and massless propagation must arise from the same completed geometry. An independently tuned optical index would undo the matter-closure result.

    \item \textbf{Normalization failure.} The single gravitational normalization left by closure must reproduce clocks, Newtonian motion, lensing, and post-Newtonian observables without sector-specific rescaling.

    \item \textbf{Empirical-distinctness problem.} Within the equilibrium minimal class the canonical closure reproduces Einstein dynamics. Any genuinely distinct phenomenology must therefore be tied to a controlled additional degree of freedom, deformed progression algebra, or non-equilibrium sector rather than reinterpretation alone.
\end{enumerate}

\section{Discussion}

The canonical calculations change the status of the progression program in several important ways.

First, the progression principle, an operational clock comparison, and the ADM lapse must not be conflated. The invariant content lies in specified clock worldlines, their comparison protocol, and relational phase/proper-time observables. A representative $\Pi_{\mathcal C}=d\tau/d\lambda$ is protocol- and parameter-dependent. The canonical lapse $N$ remains an arbitrary smearing/refoliation variable; only in an adapted foliation may one choose $N$ to represent the same clock-rate assignment. The $q=1$ statement is therefore limited pregeometrically to one temporal rate/measure factor; it is not a four-current transformation law. After geometry emerges, that temporal factor is represented inside $\sqrt{-g}=N\sqrt h$ rather than as a second physical scalar factor.

A second conceptual role is now assigned to integrability rather than to conservation alone. Quantum clock phase says what local progression means operationally, but it does not by itself determine how independently specified local advances compose. The Tomonaga--Schwinger and Dirac--Schwinger traditions show that integrability of local evolution is a genuine consistency requirement in relativistic quantum dynamics, while the Moncrief--Teitelboim result ties such integrability directly to the momentum sector of canonical GR. The present paper therefore interprets Eq.~\eqref{eq:path_independence_postulate} as a pregeometric generalization of local-evolution integrability on the quotient by spatial gauge. This is still an assumption: existing Tomonaga--Schwinger constructions already presuppose the hypersurface and causal structures that the progression program seeks to recover.

Third, the spatial completion is no longer only a suggested analogy with hypersurface geometry. The kinematic progression lemma produces a state-dependent contravariant tensor $\beta^{ij}$; Jacobi closure forces its symmetry; positivity and nondegeneracy permit $\beta^{ij}=h^{ij}$; and pointwise progression locality permits the intrinsic response to be written $\delta_Nh_{ij}=2NK_{ij}$. The canonical HDA calculation then determines that response through Eq.~\eqref{eq:K_pi_relation}. Thus the formerly free $K_{ij}$ becomes the canonical rate of spatial calibration once the minimal Hamiltonian representation is imposed.

Fourth, the nonlinear geometric dynamics are selected more strongly than in the previous version of the paper. Starting from the general $(a,b,c,d)$ constraint within the stated ultralocal quadratic class, exact HDA closure forces the DeWitt trace coefficient and the relative curvature normalization. Enlarging the spatial potential through four derivatives then removes the tested $R^2$, $R_{ij}R^{ij}$, and $D^2R$ terms. This does not prove that all conceivable progression dynamics reduce to general relativity: the result is conditional on the minimal phase space, locality, quadratic ultralocal momenta, and undeformed refoliation symmetry. It does show, however, that the Einstein spatial potential is selected rather than imported within that explicit class.

Fifth, the matter calculations strongly constrain universality across the tested principal field types, but their levels of completeness should be distinguished. Scalar and Yang--Mills closure are calculated through the relevant constraints, while the Dirac calculation fixes only the principal/characteristic co-metric; the lower-order Lorentz/gauge completion is imported from the known standard tetrad--spinor branch rather than rederived for arbitrary $\chi_D^{ij}$. Scalar self-interactions, gauge couplings, and fermion masses remain free at the level of the structure function. Linearity in the arbitrary lapse smearing provides a downstream consistency check against nonlinear $N^q$ smearing; species-dependent lapse normalizations require separate operational fixing before they can be constrained physically. The open universality problem therefore includes both a complete independent Dirac closure calculation and tests of composite, bound-state, nonminimal, and higher-derivative matter.

Sixth, the older weak-field progression results acquire a more disciplined interpretation. The static $A=\ln\Pi$, $B$ reduction is now downstream of the canonically selected ADM action. Its reduced static Dirichlet functional $c^4|\nabla A|^2/(8\pi G)$ is not an independent scalar kinetic postulate or a unique generally covariant local gravitational-energy density; $B=-A$ at leading order is the scalar-isotropic expression of the common geometric response; and the old $p=2$ optical relation follows from the same temporal and spatial geometry rather than from an independently adjustable constitutive index.

The quantum phase mechanism remains operationally useful but should not be treated as an additional force. Once universal matter closure supplies the spacetime metric, clock transitions and massive matter waves accumulate phase according to the same proper time. The weak Hamiltonian $mc^2A$ and the resulting mass-independent Ehrenfest acceleration are therefore wave-mechanical manifestations of the universal geometric coupling.

The principal unresolved problem has consequently moved. It is no longer simply to guess the dynamics of $K_{ij}$ or to test whether elementary gauge and Dirac fields can share the scalar geometry. Within the minimal canonical class those steps are fixed. The foundational burden is now to derive the upstream kinematics: the phase-defined clock increment must be promoted from an operational comparison to a genuinely local many-fingered evolution; relational progression integrability modulo spatial gauge, single-gradient leading locality, positive nondegenerate spatial calibration, canonical locality, and the undeformed HDA must then follow from that microscopic progression dynamics rather than being selected because they recover general relativity. Conservation enters as the compatible rate/measure and total-balance principle once this local evolution has acquired geometric structure. Matter closure must then be stress-tested on composite, bound-state, nonminimal, and higher-derivative systems. If those steps can be established independently, the framework would provide an integrability-and-conservation explanation of why Einstein geometry and universal coupling emerge.

\section{Conclusion}

This paper develops a relational-progression route from an operational local clock rate toward spatial geometry, canonical gravitational dynamics, and common matter characteristic cones. The construction begins with the phase-defined clock increment $d\sigma=\Pi d\lambda$, treats $\lambda$ as an arbitrary comparison parameter, and adds a local-evolution hypothesis: physical progression can be specified many-fingeredly, and different orderings of the same relational evolution must agree on physical observables modulo spatial relabeling. Tomonaga--Schwinger integrability supplies an established relativistic precedent for this consistency requirement, while the pregeometric extension remains a hypothesis. Under a local single-gradient closure this yields the normal--normal commutator
\begin{equation}
    [\delta_N,\delta_M]
    =\delta_{\beta^{ij}(N\partial_jM-M\partial_jN)}.
\end{equation}
The three-normal Jacobi identity forces $\beta^{ij}$ to be symmetric and restricts its intrinsic progression response. With positive nondegeneracy and pointwise progression locality, $\beta^{ij}$ becomes a spatial co-metric $h^{ij}$ and the response takes the kinematic form $\delta_Nh_{ij}=2NK_{ij}$.

The principal advance of the present revision is that the next dynamical step is calculated explicitly. On the minimal canonical phase space $(h_{ij},\pi^{ij})$, demanding the undeformed progression HDA for the quadratic constraint general within the stated canonical class fixes
\begin{equation}
    b=-\frac a2,
    \qquad ac=-1,
\end{equation}
leaving only the overall gravitational normalization and a cosmological term. The same calculation determines
\begin{equation}
    K_{ij}=\frac{a}{\sqrt h}
    \left(\pi_{ij}-\frac12\pi h_{ij}\right)
\end{equation}
and reconstructs the ADM action. Relaxing the spatial potential through the parity-even four-derivative basis eliminates $R^2$, $R_{ij}R^{ij}$, and $D^2R$ under exact closure. Within the stated minimal canonical class, Einstein geometrodynamics is therefore selected once the undeformed HDA has been obtained from the relational-integrability and locality assumptions. The novelty claim is not the familiar HDA-to-GR selection itself, but the proposed progression-based physical origin of that algebra.

Parallel matter calculations give a broader universality result. Scalar closure requires
\begin{equation}
    \mathsf A^{CA}\mathsf C^{ij}_{AB}=\delta^C{}_B h^{ij},
\end{equation}
Maxwell/Yang--Mills closure requires
\begin{equation}
    \alpha_E\alpha_B\chi^{ij}=h^{ij}
\end{equation}
modulo the Gauss constraint, and Dirac \emph{principal/characteristic} matching requires
\begin{equation}
    v^2\chi_D^{ij}=h^{ij}.
\end{equation}
The standard tetrad--spinor branch is known to complete the lower-order Lorentz/gauge constraints, but that full completion is not rederived here for an arbitrary trial $\chi_D^{ij}$. Thus the tested scalar, gauge, and fermionic sectors share one principal characteristic spatial co-metric, with different levels of closure completeness explicitly distinguished. Scalar potentials, gauge couplings, and fermion masses remain free because they do not alter the HDA structure function. The arbitrariness and linearity of the lapse generator select one common canonical smearing structure; possible species-dependent normalizations are not promoted here to independent physical progression charges without an operational normalization; an operational progression representative may be identified with that lapse only after choosing an adapted foliation/protocol. The scalar action reconstructs the minimally metric-coupled spacetime form explicitly, while the gauge and fermionic constraints supply the corresponding internal redundancies. The observer-current argument and the $q=1$ rate-conversion picture then provide complementary physical interpretations of the same covariant structure.

The static isotropic sector supplies a stringent consistency check. The canonically selected action reduces to the coupled $A=\ln\Pi$, $B$ system, yields $B=-A$ and the Poisson equation at weak field, produces the quadratic progression stiffness $c^4|\nabla A|^2/(8\pi G)$, recovers isotropic Schwarzschild geometry in spherical vacuum, and gives $n_{\rm eff}\simeq\Pi^{-2}$ at leading optical order. Quantum clock and matter-wave phases then respond to the same proper-time structure, with weak-field Ehrenfest motion giving $\mathbf a=-c^2\nabla\ln\Pi$.

The resulting claim is conditional and is intentionally stated narrowly. Neither the one-factor temporal rate/measure conversion nor standard Tomonaga--Schwinger theory has been shown to imply general relativity from pregeometric premises. What has been shown is that operational relational progression, together with the explicit local-evolution, integrability-modulo-gauge, locality, positivity, linear-extension, and canonical assumptions, yields the HDA candidate and strongly fixes the gravitational kinetic term, spatial curvature potential, nonlinear response, scalar/gauge closure, and Dirac principal characteristic coupling within the tested class. The decisive remaining work is to derive the many-fingered progression dynamics and its integrability principle without presupposing spacetime, establish the positive co-metric independently, and determine whether composite, nonminimal, or higher-derivative matter preserves the same closure. If those steps succeed, quantum phase would supply the operational notion of progression, integrability would explain why local progression requires geometry, canonical closure would select its minimal dynamics, and conservation would provide the balance law carried by that geometry.

\appendix
\section{Notation and interpretation}

\begin{center}
\small
\begin{tabular}{@{}>{\raggedright\arraybackslash}p{0.19\textwidth}p{0.70\textwidth}@{}}
\toprule
Symbol & Interpretation \\
\midrule
$\lambda$ & arbitrary monotonic comparison parameter; it may be used as a coordinate parameter in an adapted calculation \\
$\sigma$ & pregeometric operational clock increment defined by relative clock phase; identified with $\tau$ only after geometric completion \\
$\tau$ & ordinary metric proper time in the completed relativistic geometry; downstream identification $d\sigma=d\tau$ for an ideal clock \\
$\Pi$ & protocol- and parameter-dependent progression-rate representative, $d\sigma/d\lambda$, after a clock-comparison prescription is fixed; not progression itself \\
$\Pi_{\rm rel}$ & reparametrization-invariant ratio for specified relational clock histories under a stated comparison protocol \\
$A$ & logarithmic progression representative, $A=\ln\Pi$ in an adapted representation \\
$\delta_N$ & infinitesimal normal/update generator smeared by an arbitrary scalar $N(x)$; in an adapted operational representation a particular $N$ may encode the chosen progression-rate assignment \\
$\delta_\xi$ & infinitesimal spatial relabeling generated by $\xi^i$ \\
$\beta^{ij}$ & contravariant spatial tensor fixed kinematically by the leading path-independence commutator; Jacobi closure enforces symmetry under the stated generator-independence assumptions \\
$C^{ij}$ & free symmetric coefficient in the minimal intrinsic response $\delta_N\beta^{ij}=NC^{ij}$ \\
$C^{(ijk)}$ & Jacobi-allowed fully symmetric coefficient multiplying $\partial_kN$; removed only by pointwise progression locality or minimal tensor content \\
$h_{ij}$ & spatial metric on a local foliation, with $h^{ij}=\beta^{ij}$ after nondegeneracy and positivity are established \\
$K_{ij}$ & symmetric extrinsic-response tensor defined by $\delta_Nh_{ij}=2NK_{ij}$ in the minimal pointwise-local closure \\
$\pi^{ij}$ & canonical momentum conjugate to $h_{ij}$; HDA closure gives $\pi^{ij}=(\sqrt h/a)(K^{ij}-Kh^{ij})$ in the minimal canonical class \\
$a,b,c,d$ & coefficients of the trial gravitational Hamiltonian constraint; closure fixes $b=-a/2$ and $ac=-1$, while $d$ becomes the cosmological term \\
$\mathsf A^{AB}$ & temporal kinetic matrix in the minimal scalar matter Hamiltonian \\
$\mathsf C^{ij}_{AB}$ & spatial gradient tensor in the scalar matter Hamiltonian; closure gives $\mathsf C^{ij}_{AB}=(\mathsf A^{-1})_{AB}h^{ij}$ \\
$A_i^a,E^{ia}$ & Yang--Mills connection and conjugate electric momentum; Maxwell is the Abelian special case \\
$\chi^{ij}$ & trial gauge-field co-metric; closure fixes the physical combination $\alpha_E\alpha_B\chi^{ij}=h^{ij}$ \\
$G[\Lambda]$ & internal Yang--Mills/Maxwell Gauss generator; its appearance in the normal--normal bracket is gauge redundancy, not a second spacetime structure function \\
$\Psi$ & reduced spatial half-density Dirac spinor after elimination of the fermionic second-class constraints \\
$\chi_D^{ij}$ & candidate Dirac characteristic co-metric defined by $\{\alpha^i,\alpha^j\}=2\chi_D^{ij}$; principal/characteristic matching gives $v^2\chi_D^{ij}=h^{ij}$ \\
$J_{\rm Lor}$ & local Lorentz constraint in the tetrad--spinor canonical description \\
$B$ & isotropic spatial response in the static two-function sector \\
$S$ & nonlinear spatial-response variable, $S=A+B$ \\
$T_{\mu\nu}$ & matter stress--energy tensor obtained from the local matter action \\
$J^\mu[u]$ & energy current measured by a local observer, $-T^{\mu\nu}u_\nu$ \\
$\Phi$ & Newtonian potential, $A\simeq\Phi/c^2$ in the weak field \\
\bottomrule
\end{tabular}
\end{center}

\section{Claim hierarchy}

For clarity, the principal statements of the paper can be separated by status.

\begin{longtable}{@{}p{0.27\textwidth}p{0.63\textwidth}@{}}
\toprule
Status & Statement \\
\midrule
\endfirsthead
\toprule
Status & Statement \\
\midrule
\endhead
\midrule
\multicolumn{2}{r}{\small Continued on next page}\\
\endfoot
\bottomrule
\endlastfoot
Operational definition & Clock-transition phase defines the pregeometric increment $d\sigma=-(\hbar/\Delta E)d(\Delta\theta)$; proper time is identified only after geometric completion. $\lambda$ itself need not be physical. \\
Rate/measure-conversion result & Expressing a local rate per unit $\lambda$ introduces one temporal factor of $\Pi$. This does not by itself transform a full pregeometric four-current or establish a four-dimensional divergence law. \\
Integrability-motivation statement & Tomonaga--Schwinger/Dirac--Schwinger theory provides an established precedent that independently specified local relativistic evolution must satisfy integrability conditions. The manuscript elevates this to a pregeometric hypothesis of relational progression integrability modulo spatial relabeling; the precedent motivates but does not derive that hypothesis. \\
Kinematic lemma & On an assumed smooth three-dimensional relational state manifold, under real linear extension, vertical-curvature integrability modulo labels, local single-gradient closure, spatial covariance, and normal/tangential generator independence, $[\delta_N,\delta_M]=\delta_{\beta^{ij}(N\partial_jM-M\partial_jN)}$. \\
Jacobi compatibility result & Independence of normal and tangential generators plus the three-normal Jacobi identity gives $\beta^{ij}=\beta^{ji}$ and, at first-gradient order, $\delta_N\beta^{ij}=NC^{(ij)}+C^{(ijk)}\partial_kN$. \\
Conditional co-metric response & Positivity/nondegeneracy plus the universal identification of the commutator tensor with the directional-transport calibration identify $\beta^{ij}=h^{ij}$; pointwise progression locality removes $C^{(ijk)}\partial_kN$ and permits $\delta_Nh_{ij}=2NK_{ij}$. \\
Canonical gravity closure result & On the minimal local $(h_{ij},\pi^{ij})$ phase space with quadratic ultralocal momenta and potential $cR+d$, exact progression-HDA closure gives $b=-a/2$, $ac=-1$, and $K_{ij}=(a/\sqrt h)(\pi_{ij}-\pi h_{ij}/2)$. \\
Four-derivative potential result & Within the parity-even basis $d+cR+\alpha R^2+\zeta R_{ij}R^{ij}+\gamma D^2R$, undeformed HDA closure requires $\alpha=\zeta=\gamma=0$. This does not exclude enlarged phase spaces or deformed algebras. \\
Scalar matter closure & For the separable scalar Hamiltonian in Eq.~\eqref{eq:general_matter_H}, closure requires $\mathsf A^{CA}\mathsf C^{ij}_{AB}=\delta^C{}_B h^{ij}$ and reconstructs one minimally coupled spacetime metric. \\
Gauge-field closure result & For the trial Maxwell/Yang--Mills Hamiltonian in Eq.~\eqref{eq:YM_trial_H}, closure modulo the Gauss constraint requires $\alpha_E\alpha_B\chi^{ij}=h^{ij}$. Gauge coupling strength remains free. \\
Dirac principal/characteristic-matching result & On the reduced trial spinor phase space, the calculated principal normal--normal bracket has structure function $v^2\chi_D^{ij}$; characteristic matching requires $v^2\chi_D^{ij}=h^{ij}$ and leaves the fermion mass free. Full lower-order closure for arbitrary trial $\chi_D^{ij}$ is not shown; the standard tetrad--spinor branch is known to close modulo local Lorentz and gauge constraints. \\
Conditional smearing/lapse result & Physical forward-progression profiles are positive; the canonical algebra additionally assumes a consistent real linear extension to arbitrary smearings $N$. Only after closure may an adapted foliation/protocol choose a particular $N$ to represent an operational progression assignment. \\
Covariant consequence & The reconstructed matter action has $\sqrt{-g}=N\sqrt h$ and yields $\nabla_\mu T^{\mu\nu}=0$ on shell; the observer-current argument supplies a complementary pointwise interpretation. \\
Static reduction result & The canonically selected action yields the weak-field quadratic progression stiffness, $B=-A$ at leading order, exact isotropic Schwarzschild vacuum geometry, and $n_{\rm eff}\simeq\Pi^{-2}$ at leading optical order. \\
Open foundational step & Derive the assumed smooth relational state manifold, many-fingered local progression dynamics, vertical-curvature integrability modulo spatial gauge, and normal/tangential generator independence without presupposing spacetime; then derive single-gradient kinematics, positive nondegenerate spatial calibration, the minimal canonical phase space, and the undeformed HDA rather than selecting them as the closure class. This is the central unresolved bridge. \\
Open matter step & Test composite and bound-state matter, nonminimal curvature couplings, higher-derivative effective interactions, and any enlarged physical phase space for preservation of the same physical co-metric and lapse. \\
Open new-physics step & If distinct phenomenology is intended, identify a controlled additional degree of freedom, deformed algebra, or non-equilibrium sector after the equilibrium minimal theory is fixed. \\
\end{longtable}

\begin{thebibliography}{99}

\bibitem{Stodolsky1979}
L. Stodolsky,
``Matter and light wave interferometry in gravitational fields,''
\emph{General Relativity and Gravitation} \textbf{11}, 391--405 (1979).
\href{https://doi.org/10.1007/BF00759302}{doi:10.1007/BF00759302}.

\bibitem{Dimopoulos2008}
S. Dimopoulos, P. W. Graham, J. M. Hogan, and M. A. Kasevich,
``General relativistic effects in atom interferometry,''
\emph{Physical Review D} \textbf{78}, 042003 (2008).
\href{https://doi.org/10.1103/PhysRevD.78.042003}{doi:10.1103/PhysRevD.78.042003}.

\bibitem{Zych2011}
M. Zych, F. Costa, I. Pikovski, and \v{C}. Brukner,
``Quantum interferometric visibility as a witness of general relativistic proper time,''
\emph{Nature Communications} \textbf{2}, 505 (2011).
\href{https://doi.org/10.1038/ncomms1498}{doi:10.1038/ncomms1498}.

\bibitem{Roura2020}
A. Roura,
``Gravitational redshift in quantum-clock interferometry,''
\emph{Physical Review X} \textbf{10}, 021014 (2020).
\href{https://doi.org/10.1103/PhysRevX.10.021014}{doi:10.1103/PhysRevX.10.021014}.

\bibitem{Ufrecht2020}
C. Ufrecht, F. Di Pumpo, A. Friedrich, A. Roura, C. Schubert, D. Schlippert,
E. M. Rasel, W. P. Schleich, and E. Giese,
``Atom-interferometric test of the universality of gravitational redshift and free fall,''
\emph{Physical Review Research} \textbf{2}, 043240 (2020).
\href{https://doi.org/10.1103/PhysRevResearch.2.043240}{doi:10.1103/PhysRevResearch.2.043240}.

\bibitem{Roura2021}
A. Roura, C. Schubert, D. Schlippert, and E. M. Rasel,
``Measuring gravitational time dilation with delocalized quantum superpositions,''
\emph{Physical Review D} \textbf{104}, 084001 (2021).
\href{https://doi.org/10.1103/PhysRevD.104.084001}{doi:10.1103/PhysRevD.104.084001}.

\bibitem{Tomonaga1946}
S.-I. Tomonaga,
``On a relativistically invariant formulation of the quantum theory of wave fields,''
\emph{Progress of Theoretical Physics} \textbf{1}, 27--42 (1946).
\href{https://doi.org/10.1143/PTP.1.27}{doi:10.1143/PTP.1.27}.

\bibitem{Schwinger1948}
J. Schwinger,
``Quantum electrodynamics. I. A covariant formulation,''
\emph{Physical Review} \textbf{74}, 1439--1461 (1948).
\href{https://doi.org/10.1103/PhysRev.74.1439}{doi:10.1103/PhysRev.74.1439}.

\bibitem{Schwinger1963}
J. Schwinger,
``Energy and momentum density in field theory,''
\emph{Physical Review} \textbf{130}, 800 (1963).
\href{https://doi.org/10.1103/PhysRev.130.800}{doi:10.1103/PhysRev.130.800}.

\bibitem{Fukuyama1973}
T. Fukuyama,
``Integrability condition and energy-density commutation relation,''
\emph{Progress of Theoretical Physics} \textbf{49}, 1340--1351 (1973).
\href{https://doi.org/10.1143/PTP.49.1340}{doi:10.1143/PTP.49.1340}.

\bibitem{MoncriefTeitelboim1972}
V. Moncrief and C. Teitelboim,
``Momentum constraints as integrability conditions for the Hamiltonian constraint in general relativity,''
\emph{Physical Review D} \textbf{6}, 966--969 (1972).
\href{https://doi.org/10.1103/PhysRevD.6.966}{doi:10.1103/PhysRevD.6.966}.

\bibitem{Hoehn2024}
P. A. H\"ohn, A. Russo, and A. R. H. Smith,
``Matter relative to quantum hypersurfaces,''
\emph{Physical Review D} \textbf{109}, 105011 (2024).
\href{https://doi.org/10.1103/PhysRevD.109.105011}{doi:10.1103/PhysRevD.109.105011};
\href{https://arxiv.org/abs/2308.12912}{arXiv:2308.12912}.

\bibitem{Barbour2002}
J. Barbour, B. Z. Foster, and N. {\'O} Murchadha,
``Relativity without relativity,''
\emph{Classical and Quantum Gravity} \textbf{19}, 3217 (2002).
\href{https://doi.org/10.1088/0264-9381/19/12/308}{doi:10.1088/0264-9381/19/12/308};
\href{https://arxiv.org/abs/gr-qc/0012089}{arXiv:gr-qc/0012089}.

\bibitem{AndersonBarbour2002}
E. Anderson and J. Barbour,
``Interacting vector fields in relativity without relativity,''
\emph{Classical and Quantum Gravity} \textbf{19}, 3249 (2002).
\href{https://doi.org/10.1088/0264-9381/19/12/309}{doi:10.1088/0264-9381/19/12/309};
\href{https://arxiv.org/abs/gr-qc/0201092}{arXiv:gr-qc/0201092}.

\bibitem{Teitelboim1973}
C. Teitelboim,
``How commutators of constraints reflect the spacetime structure,''
\emph{Annals of Physics} \textbf{79}, 542--557 (1973).
\href{https://doi.org/10.1016/0003-4916(73)90096-1}{doi:10.1016/0003-4916(73)90096-1}.

\bibitem{HKT1976}
S. A. Hojman, K. Kucha\v{r}, and C. Teitelboim,
``Geometrodynamics regained,''
\emph{Annals of Physics} \textbf{96}, 88--135 (1976).
\href{https://doi.org/10.1016/0003-4916(76)90112-3}{doi:10.1016/0003-4916(76)90112-3}.

\bibitem{Duell2018}
M. D\"ull, F. P. Schuller, N. Stritzelberger, and F. Wolz,
``Gravitational closure of matter field equations,''
\emph{Physical Review D} \textbf{97}, 084036 (2018).
\href{https://doi.org/10.1103/PhysRevD.97.084036}{doi:10.1103/PhysRevD.97.084036};
\href{https://arxiv.org/abs/1611.08878}{arXiv:1611.08878}.

\bibitem{Thiemann2024}
T. Thiemann,
``Non-degenerate metrics, hypersurface deformation algebra, non-anomalous representations and density weights in quantum gravity,''
\emph{General Relativity and Gravitation} \textbf{56}, 122 (2024).
\href{https://doi.org/10.1007/s10714-024-03313-w}{doi:10.1007/s10714-024-03313-w}.

\bibitem{Bojowald2025}
M. Bojowald, E. I. Duque, and A. Shah,
``Hypersurface deformations: Off-shell properties on phase space,''
\emph{Physical Review D} \textbf{111}, 124048 (2025).
\href{https://doi.org/10.1103/c1kn-xljx}{doi:10.1103/c1kn-xljx};
\href{https://arxiv.org/abs/2410.18807}{arXiv:2410.18807}.

\bibitem{KucharStone1987}
K. V. Kucha\v{r} and C. L. Stone,
``A canonical representation of space-time diffeomorphisms for the parametrized Maxwell field,''
\emph{Classical and Quantum Gravity} \textbf{4}, 319--328 (1987).
\href{https://doi.org/10.1088/0264-9381/4/2/013}{doi:10.1088/0264-9381/4/2/013}.

\bibitem{NelsonTeitelboim1978}
J. E. Nelson and C. Teitelboim,
``Hamiltonian formulation of the theory of interacting gravitational and electron fields,''
\emph{Annals of Physics} \textbf{116}, 86--104 (1978).
\href{https://doi.org/10.1016/0003-4916(78)90005-2}{doi:10.1016/0003-4916(78)90005-2}.

\bibitem{Deser1970}
S. Deser,
``Self-interaction and gauge invariance,''
\emph{General Relativity and Gravitation} \textbf{1}, 9--18 (1970).
\href{https://doi.org/10.1007/BF00759198}{doi:10.1007/BF00759198};
\href{https://arxiv.org/abs/gr-qc/0411023}{arXiv:gr-qc/0411023}.

\bibitem{Deser2010}
S. Deser,
``Gravity from self-interaction redux,''
\emph{General Relativity and Gravitation} \textbf{42}, 641--646 (2010).
\href{https://doi.org/10.1007/s10714-009-0915-7}{doi:10.1007/s10714-009-0915-7};
\href{https://arxiv.org/abs/0910.2975}{arXiv:0910.2975}.

\bibitem{Lovelock1971}
D. Lovelock,
``The Einstein tensor and its generalizations,''
\emph{Journal of Mathematical Physics} \textbf{12}, 498--501 (1971).
\href{https://doi.org/10.1063/1.1665613}{doi:10.1063/1.1665613}.

\bibitem{Jacobson1995}
T. Jacobson,
``Thermodynamics of spacetime: The Einstein equation of state,''
\emph{Physical Review Letters} \textbf{75}, 1260--1263 (1995).
\href{https://doi.org/10.1103/PhysRevLett.75.1260}{doi:10.1103/PhysRevLett.75.1260};
\href{https://arxiv.org/abs/gr-qc/9504004}{arXiv:gr-qc/9504004}.

\bibitem{Eling2006}
C. Eling, R. Guedens, and T. Jacobson,
``Non-equilibrium thermodynamics of spacetime,''
\emph{Physical Review Letters} \textbf{96}, 121301 (2006).
\href{https://doi.org/10.1103/PhysRevLett.96.121301}{doi:10.1103/PhysRevLett.96.121301};
\href{https://arxiv.org/abs/gr-qc/0602001}{arXiv:gr-qc/0602001}.

\bibitem{Jacobson2016}
T. Jacobson,
``Entanglement equilibrium and the Einstein equation,''
\emph{Physical Review Letters} \textbf{116}, 201101 (2016).
\href{https://doi.org/10.1103/PhysRevLett.116.201101}{doi:10.1103/PhysRevLett.116.201101};
\href{https://arxiv.org/abs/1505.04753}{arXiv:1505.04753}.

\bibitem{JacobsonVisser2019}
T. Jacobson and M. R. Visser,
``Spacetime equilibrium at negative temperature and the attraction of gravity,''
\emph{International Journal of Modern Physics D} \textbf{28}, 1944016 (2019).
\href{https://doi.org/10.1142/S0218271819440167}{doi:10.1142/S0218271819440167};
\href{https://arxiv.org/abs/1904.04843}{arXiv:1904.04843}.

\end{thebibliography}

\end{document}
